We show consistency of our proposed optimal ITR estimator \( \doptHAT(Z) \), as well as the random survival forest in the single stage setting and our proposed RCIF. \if0\prelim The proofs of all theorems are presented in the Supplementary Material. \fi We use the definition of regular trees and random-split trees from the random forest literature \citep{meinshausen_quantile_2006, wager_adaptive_2016, cui_consistency_2022}.
\begin{definition}[Random-split and \( \alpha \)-regular trees and forests]
\label{def:regsplotalphatree}
A tree is called a random-split tree if each feature is given a minimum probability \( (\psi/d) \) of being the split variable at each intermediate node, where \( 0 < \psi < 1 \) and \( d \) is the dimension of the feature space. A tree is \( \alpha \)-regular if every daughter node has at least an \( \alpha \) fraction of the training sample in the corresponding parent node. A random forest is called a random-split (\( \alpha \)-regular) forest if each of its member trees is random-split (\( \alpha \)-regular).
\end{definition}

In addition to Assumptions \ref{assump:Causal1}-\ref{assump:Causal3} in Section \ref{sec:proposedmethod}, we also assume the following assumptions.
% survival assumption
Assumption \ref{assump:CuiAssump1} is important for asymptotic results, and is typical for survival analysis \citep{fleming_counting_2011,cui_tree_2017}. This assumption ensures we only calculate each survival or cumulative incidence curve up to a time point with sufficiently large sample size.

\begin{assumption}
\label{assump:CuiAssump1}
There exists a fixed maximum follow-up time $0 < \tau < \infty$ and a constant $\tauUB \in (0, 1)$, such that $Pr(X_\tau \geq \tau | Z = z) \geq \tauUB$, for all $z \in \mathcal{Z}$.
\end{assumption}

% INPUTTING ASSUMPTIONS for forest consistency
% Assumptions for Forest Consistency
Below are assumptions used for consistency of both forests. Assumptions \ref{assump:Cho1}-\ref{assump:Cho3} are necessary and realistic for the tree-based estimator terminal node size and splitting rules. Assumption \ref{assump:Cho1} requires the terminal node size to be not only sufficiently large absolutely but also sufficiently small relative to the sample size. Assumption \ref{assump:Cho2} follows from Definition \ref{def:regsplotalphatree}. Assumption \ref{assump:Cho3} adds smoothness to the survival function and censoring cumulative distribution function. Assumption \ref{assump:CuiAssump2} follows from \cite{cui_tree_2017} and \cite{wager_adaptive_2016}, which allows for dependency between covariates. Thus, non-hyper-rectangular or categorical history spaces are allowed with this assumption. Assumption \ref{khbddtv} is necessary for the random cumulative incidence forest because there is non-differentiability of the treatment policy.

\begin{assumption}[Terminal node size, polynomial]
\label{assump:Cho1}
The minimum size \( n_{\text{min}} \) of the terminal nodes grows at the rate \( n_{\text{min}} \bowtie n^\beta \), \( \frac{1}{2} < \beta < 1 \), where \( a \bowtie b \) implies that both \( a = O(b) \) and \( b = O(a) \).
\end{assumption}

\begin{assumption}
[\( \alpha \)-regular and random-split trees]
\label{assump:Cho2}
Trees are \( \alpha \)-regular and random-split with a constant \( 0 < \psi < 1 \).
\end{assumption}

\begin{assumption}[Lipschitz continuous and bounded survival and censoring cumulative densities]
\label{assump:Cho3}
There exist constants \( \lipS \) and \( \lipG \) such that \( |\survfunc(t | h_1) - \survfunc(t | h_2)| \leq \lipS\|h_1 - h_2\| \) and \( |G(t | h_1) - G(t | h_2)| \leq \lipG \|h_1 - h_2\| \) for all \( h_1, h_2 \in \history \), \( t \in [0, \tau] \), where \( G \) is the censoring cumulative distribution function (CDF) and \( \survfunc(\tau^- \mid \histlower)G(\tau^- \mid \histlower) > c_1 \) for all \( \histlower \) and some constant \( c_1 > 0 \).
\end{assumption}

\begin{assumption}
    \label{assump:CuiAssump2}
    The covariates $Z \in [0; 1]^d$ are distributed according to a density function $p(\cdot)$ satisfying $1/\assumpcovbound \leq p(z) \leq \assumpcovbound$, for all $z$ and some $\assumpcovbound \geq 1$.
\end{assumption}

\begin{assumption}
\label{khbddtv}
The kernel \( \kernel \) for the random cumulative incidence forest has bounded total variation. Specifically, there exists a constant \( C \) such that the total variation of \( \kernel \) satisfies
\[
\int_{-\infty}^{\infty} \abs{\kernel^\prime(x)} \, dx \leq C,
\]
where \( \kernel^\prime(x) \) denotes the derivative of the kernel \( \kernel \).
\end{assumption}

% ITR consistency assumptions
% Assumptions for ITR Consistency
Below are assumptions for the proposed estimator for consistency. We assume bounding in Assumptions \ref{assump:boundeddensity}-\ref{assump:boundedV0} to account for nonregularity of the value function. This is a necessary and standard assumption for non-regularity conditions.

\begin{assumption}
\label{assump:boundeddensity}
Define the random function $B(a,Z) = \valueoneTRUE - (1-\tol)\valueonemaxTRUE$, $a \in \mathcal{A}$. Assume $B(a,Z)$ has continuous, bounded densities at 0, for each $a \in \mathcal{A}$.
\end{assumption}

\begin{assumption}
\label{assump:boundedV0}
    $\sup_{a,Z} \left\| \valueTRUE \right\|$ is bounded by some $\sevenbound < \infty$.
\end{assumption}

%%%% ITR
\subsection{Individualized Treatment Rule Estimator}
\label{Sec:ITRestimator_Consistency}

For any fixed value of covariate $Z$, the difference in value functions between true optimal ITR and estimated optimal ITR based on Algorithm \ref{alg:Overview} is
\begin{equation}
\label{eqn:ValueDiff}
    \mathbb{V}(\doptTRUE(Z)) - \mathbb{V}(\doptHAT(Z)) = \sum_{a \in \mathcal{A}} \mathbb{E}\left[ \valueTRUE
\cdot \left( \mathbf{1}\{\doptTRUE(Z) = a\} - \mathbf{1}\{\doptHAT(Z) = a\} \right)
\right]
\end{equation}
where \( \mathbb{V}(d(Z)) \) is $\phi$ of a treatment regime \( d(Z) \) if all individuals in a population follow the treatment regime. Showing Equation \cref{eqn:ValueDiff} converges in probability to $0$ proves that $\doptHAT$ converges in probability to the true $\doptTRUE$. The following theorem states that the individualized treatment regime estimator has a value consistent for the value of the optimal regime.

% ITR estimator CONSISTENCY Theorem and Proof
%%%%%%%%%%%%%%%%%%%%%%%%%%%%%%%%%%%%%%%%%%%%%%%%%%%%%%%%%%%
%%%%%%%%%%%%%%%%%%%%%%%%%%%%%%%%%%%%%%%%%%%%%%%%%%%%%%%%%%%
%%%%%%%%%%%%%%%%%%%%%%%%%%%%%%%%%%%%%%%%%%%%%%%%%%%%%%%%%%%
% Consistency of ITR estimator
\newcommand{\ITRconsistency}{
\begin{theorem} (Consistency of the individualized treatment regime estimator) 
\label{thm:ITR_Consistency}
Assume Assumptions \ref{assump:CuiAssump1}-\ref{assump:CuiAssump2} hold. Let $\doptHAT$ denote the optimal ITR estimator that is built following Algorithm \ref{alg:Overview} and Assumptions \ref{assump:boundeddensity}-\ref{assump:boundedV0}. Assume Assumptions \ref{assump:Causal1}-\ref{assump:Causal3} hold and that the random forests are built for each treatment arm. Then, for %each $\phi = \phi^{\mu}, \phi^{\sigma,t_0}$ where 
$\phi^\mu = \begin{pmatrix}
    \phi^\mu_1 \\ \phi^\mu_2
\end{pmatrix},$ % and $\phi^{\sigma, t_0} = \begin{pmatrix} \phi^{\sigma, t_0}_1 \\ \phi^{\sigma, t_0}_2 \end{pmatrix},$
\[ \lim_{n \to \infty} \| \mathbb{V}(\doptTRUE(Z)) - \mathbb{V}(\doptHAT(Z)) \| \xrightarrow{p} 0\]
where $\|\cdot\|$ is defined as the $L_2$ norm.
\end{theorem}
}

\newcommand{\ITRconsistencyPROOF}{
\begin{proof}[Proof of Theorem \ifnum\jasa=1 1\else \ref{thm:ITR_Consistency}\fi]
%We know that, by properties of norm, $\left\| \begin{pmatrix} a_1 \\ a_2 \end{pmatrix} + \begin{pmatrix} b_1 \\ b_2 \end{pmatrix}\right\| \leq \left\| \begin{pmatrix} a_1 \\ a_2 \end{pmatrix} \right\| + \left\| \begin{pmatrix} b_1 \\ b_2 \end{pmatrix} \right\|$ and $\left\| c\begin{pmatrix} a_1 \\ a_2 \end{pmatrix} \right\| =  |c| \left\|\begin{pmatrix} a_1 \\ a_2 \end{pmatrix}\right\|$ where $c$ is a scalar. Also, by property of expectation, we know that $\|\mathbb{E}(\cdot)\| \leq \mathbb{E}(\|\cdot\|)$.

Note that
\begin{equation*}
\begin{split}
\left\| \mathbb{V}(\doptTRUE(Z)) - \mathbb{V}(\doptHAT(Z)) \right\| 
% line 1
&= \left\|\sum_{a \in \mathcal{A}} \mathbb{E}_Z\left[ \valueTRUE
\cdot \left( \mathbf{1}\{\doptTRUE(Z) = a\} - \mathbf{1}\{\doptHAT(Z) = a\} \right)
\right] \right\| \\
% line 2
&\leq \sum_{a \in \mathcal{A}} \left\|\mathbb{E}_Z\left[ \valueTRUE
\cdot \left( \mathbf{1}\{\doptTRUE(Z) = a\} - \mathbf{1}\{\doptHAT(Z) = a\} \right)
\right] \right\| \\
% line 3
&\leq \sum_{a \in \mathcal{A}} \mathbb{E}_Z \left\| \left[ \valueTRUE
\cdot \left( \mathbf{1}\{\doptTRUE(Z) = a\} - \mathbf{1}\{\doptHAT(Z) = a\} \right)
\right] \right\| \\
% line 4
&\leq \sum_{a \in \mathcal{A}} \mathbb{E}_Z \left[ \left\| \valueTRUE \right\|
\cdot \left| \mathbf{1}\{\doptTRUE(Z) = a\} - \mathbf{1}\{\doptHAT(Z) = a\} \right|
\right] \\
% line 5
&= \mathbb{E}_Z \left[ \sum_{a \in \mathcal{A}} \left( \mathbf{1} \{\doptTRUE(Z) \neq \doptHAT(Z)\} \left\|\valueTRUE\right\| \right) \right] \\
% line 6
&\leq \mathbb{E}_Z \left( \mathbf{1} \{\doptTRUE(Z) \neq \doptHAT(Z)\} \right) \sup_{a,Z} \left\| \valueTRUE \right\| \\
% line 7
&= \text{Pr}(\doptTRUE(Z) \neq \doptHAT(Z) | \doptHAT) \sup_{a,Z} \left\| \valueTRUE \right\|\\
% line 8
&\leq \sevenbound \text{Pr}(\doptTRUE(Z) \neq \doptHAT(Z) | \doptHAT).\\
\end{split} % split gives one number; align gives a number per line
\end{equation*} The third to last line holds because 1) only one term in the summation will have a non-zero contribution and 2) the supremum represents the ``worst case scenario" over all fixed combinations of $a$ and $z$ and thus does not depend on the random variables. The last line holds because we assumed $\sup_{a,Z} \left\| \valueTRUE \right\|$ is bounded by $\sevenbound$ via Assumption \ifnum\jasa=1 11\else \ref{assump:boundedV0} \fi. Thus, the norm of this distance, $\mathbb{V}(\doptTRUE) - \mathbb{V}(\doptHAT)$, goes to $0$ in probability if $\text{Pr}(\doptTRUE(Z) \neq \doptHAT(Z)|\hat{d}_\text{opt}) \xrightarrow{p} 0$.\\

We complete the proof by first showing
\begin{align}
    \label{eqn:step1}
        \sum_{a \in \mathcal{A}} \left| I\{ \valueoneTRUE - (1-\tol)\valueonemaxTRUE \ge 0 \} - I\{ \valueoneHAT - (1-\tol)\valueonemaxHAT \ge 0 \} \right| \rightarrow 0.
        \end{align}
    %(Note: we take the difference because both are random (one relies on Z and the other is an estimator)).
The proof of \cref{eqn:step1} will be presented shortly. Now, \cref{eqn:step1} implies 
\begin{align}
    \label{eqn:star}
        \argmax_{a \in \mathcal{A}: \valueoneHAT \ge (1-\tol)\valueonemaxHAT} \valuetwoHAT \xrightarrow{p} \argmax_{a \in \mathcal{A}: \valueoneTRUE \ge (1-\tol)\valueonemaxTRUE} \valuetwoHAT.
    \end{align} 
Theorem \ifnum \jasa=1 3 \else \ref{thm:RCIF_Consistency} \fi in our paper states that RCIFs are consistent. %replace <- with theorem when typed up.
Thus, by the Continuous Mapping Theorem, $\valuetwoHAT \xrightarrow{p} \valuetwoTRUE$. %$\valuetwoHAT$ is bounded.
Then,
\begin{align}
\label{eqn:step2}
        \sup_{a \in \mathcal{A}} \, \mathbb{E} \, \left| \valuetwoHAT - \valuetwoTRUE \right| \xrightarrow{p} 0,
    \end{align} 
which implies that the right hand side of \cref{eqn:star} converges to \\$\argmax_{a \in \mathcal{A}: \valueoneTRUE \ge (1-\tol)\valueonemaxTRUE} \valuetwoTRUE$. Thus, since $Z$ is random, $$\text{Pr}(\doptTRUE(Z) \neq \doptHAT(Z) | \doptHAT) \xrightarrow{p} 0.$$ This completes the proof of Theorem \ifnum\jasa=1 1\else \ref{thm:ITR_Consistency}\fi.
\end{proof}
}

\newcommand{\eqnsteponePROOF}{
%\textbf{Proof of \cref{eqn:step1}:}
\begin{proof}[Proof of \cref{eqn:step1}]
Let $A_0 = \valueoneTRUE - (1-\tol)\valueonemaxTRUE$
and $\widehat{A}_n = \valueoneHAT - (1-\tol)\valueonemaxHAT$. These are functions of random variable $Z$. Then, we can re-write \cref{eqn:step1} as 
$$\sum_{a \in \mathcal{A}} \left| I\{ A_0 \ge 0 \} - I \{ \widehat{A}_n \ge 0 \} \right| \rightarrow 0.$$
\begin{equation}
\label{eqn:terms12}
\begin{split}
    \left| I\{ \widehat{A}_n \ge 0 \} - I \{A_0 \ge 0 \} \right| 
    &= I\{ \widehat{A}_n < 0 \le A_0 \} + I \{A_0 < 0 \le \widehat{A}_n \} \\
    &= I\{ \widehat{A}_n - A_0 < -A_0 \leq 0 \} + I \{0 < -A_0 \le \widehat{A}_n - A_0 \} \\
    &= I\{ 0 \le A_0 < -(\widehat{A}_n - A_0) \} + I \{- (\widehat{A}_n - A_0) \le A_0 < 0 \}. \\
\end{split}
\end{equation}
In \cref{eqn:terms12}, the first equality holds because these are the only possible indicators such that their sum is equal to $1$. We get the 2nd equality by subtracting $A_0$ on both sides within the indicator functions. The third equality is due to multiplying by $-1$ and flipping the inequalities within the indicators.

Since we know $\valueoneHAT \xrightarrow{p} \valueoneTRUE$ and $\valueonemaxHAT \xrightarrow{p} \valueonemaxTRUE$ as shown in \cite{cho_multi-stage_2023}; and $\widehat{A}_n$ and $A_0$ are linear combinations of $\valueoneHAT$ and $\valueonemaxHAT$, and $\valueoneTRUE$ and $\valueonemaxTRUE$, respectively, then %Theorem \ref{thm:consistency_lincomb},
$\widehat{A}_n \xrightarrow{p} A_0$ because the linear combination of consistent estimators is consistent. This means
\begin{align}
\label{eqn:AnA0} 
\widehat{A}_n - A_0 \xrightarrow{p} 0.
\end{align}

By definition of convergence in probability, \cref{eqn:AnA0} implies $\exists\epsilon_n: Pr(|\widehat{A}_n - A_0|>\epsilon_n) \rightarrow 0 \text{ as } \epsilon_n \rightarrow 0$ for a fixed sequence $\epsilon_n$. Then, using Assumption \ifnum\jasa=1 10 \else \ref{assump:boundeddensity} \fi and \cref{eqn:AnA0}, we can bound \cref{eqn:terms12}. Starting with the first term of \cref{eqn:terms12},
\begin{equation*}
\label{eqn:term1}
\begin{split}
&I\{ 0 \le A_0 < -(\widehat{A}_n - A_0), |\widehat{A}_n-A_0|\le \epsilon_n \}  + I\{ 0 \le A_0 < -(\widehat{A}_n - A_0), |\widehat{A}_n-A_0| > \epsilon_n \} \le \\&I\{ 0 \le A_0 < \epsilon_n\} + I \{ |\widehat{A}_n - A_0| > \epsilon_n \}, \\    
\end{split}
\end{equation*}
 where $I\{ 0 \le A_0 < \epsilon_n\} \rightarrow 0$ by Assumption \ifnum\jasa=1 10 \else \ref{assump:boundeddensity}\fi (bounded density assumption) and
$I \{ |\widehat{A}_n - A_0| > \epsilon_n \} \rightarrow 0$ by \cref{eqn:AnA0} (and definition of consistency). Specifically, assuming Assumption \ifnum\jasa=1 10\else \ref{assump:boundeddensity}\fi,
$Pr(0\le A_0 \le \epsilon_n) = \int^{\epsilon_n}_{0} f_{A_0}(a_0)da_0 \approx f_{a_0}(0)\epsilon_n \le b\epsilon_n \rightarrow 0 $ as $\epsilon_n \rightarrow 0 $ for some bound $b$. The second term can be shown to be bounded similarly. 
\begin{equation*}
\begin{split}
&I \{- (\widehat{A}_n - A_0) \le A_0 < 0 \} = \\
&I \{- (\widehat{A}_n - A_0) \le A_0 < 0 , |\widehat{A}_n - A_0| \le \epsilon_n\} +
I \{- (\widehat{A}_n - A_0) \le A_0 < 0 , |\widehat{A}_n - A_0| > \epsilon_n\} \leq\\
&I(-\epsilon_n \leq A_0 < 0) + I(|\widehat{A}_n - A_0| > \epsilon_n).
\end{split}
\end{equation*} The first term for the inequality is because $-(\widehat{A}_n - A_0) < 0$, $\widehat{A}_n - A_0$ must be positive; thus, $-\epsilon_n \leq \widehat{A}_n - A_0 \leq \epsilon_n$.  The second term is because the indicator is a subset.

Because of Equation \cref{eqn:AnA0}, $Pr(|\widehat{A}_n - A_0| > \epsilon_n) \rightarrow 0$, and $Pr(-\epsilon_n \leq A_0 < 0) \rightarrow 0$ due to Assumption \ifnum\jasa=1 10\else \ref{assump:boundeddensity}\fi.
Specifically, $I(-\epsilon_n \leq A_0 < 0)  \implies Pr(-\epsilon_n \le A_0 < 0) \approx \epsilon_n f_{a_0}(0) < \epsilon_n b \rightarrow 0$ as $\epsilon_n \rightarrow 0$.
This completes the proof of \cref{eqn:step1}.
\end{proof}
}

\begin{comment}
MARCH 2024 UPDATE: WE CAN IGNORE BELOW I THINK.

We want to use the below equation. I think he said that, as long as the density of the quantity $\valueoneTRUE - (1-\tol)\valueonemaxTRUE$ is bounded then the whole thing converges to 0?
\begin{equation}
%\begin{align}
    \lim_{n \to \infty} E\left(\left|I( \valueoneHAT - (1-\tol)\valueonemaxHAT \ge 0) - I( \valueoneTRUE - (1-\tol)\valueonemaxTRUE \ge 0 )\right|\right) = 0 \\
%\end{align}
\end{equation}
We want to show that $\valueoneHAT - (1-\tol)\valueonemaxHAT -\xrightarrow{p} \valueoneTRUE - (1-\tol)\valueonemaxTRUE$
Let $X_n \xrightarrow{p} \mu$ and $Y_n \xrightarrow{p} \gamma $. Then, $Pr(|X_n - \mu| > \epsilon \text{ or } |Y_n - \gamma|>\epsilon) \leq Pr(|X_n - \mu|>\epsilon) + Pr(|Y_n - \gamma|>\epsilon)$. If $n \rightarrow \infty$, then both $Pr(|X_n - \mu|>\epsilon) \rightarrow 0$ and $Pr(|Y_n - \gamma|>\epsilon) \rightarrow 0$. We want to use this as the generic idea for $X_n$ being OS and $Y_n$ being CIF.
\end{comment}

\begin{comment}
MARCH 2024 UPDATE: DEFINITELY IGNORE BELOW. PROOF IS COMPLETE.

\subsubsection{OLD STUFF: Notes from before 2/20/24 meeting}

Sample size $n$ such that $\hat{d}_n$. Let 
\[
V(\hat{d}_n) = \begin{pmatrix}
V_1(\hat{d}_n) \\ V_2(\hat{d}_n)
\end{pmatrix}
\]

We want to show:
\[
\lim_{{n \to \infty}} \bigg( \frac{V_1(d_*) - V_1(\hat{d}_n)}{V_1(d_*)} - \tol \bigg)_+ \xrightarrow{p} 0
\]
\\

Our optimal $d_{**}$ satisfies the following two criteria:
\begin{enumerate}
    \item $E\left[ I(\frac{V_1(d_*|Z)-V_1(d_{**}|Z)}{V_1(d_*|Z)} \leq \tol ) \right] = 1$
    \item maximize $E\left[V_2(d_{**}|Z) \;\middle|\; I(\frac{V_1(d_*|Z)-V_1(d_{**}|Z)}{V_1(d_*|Z)} \leq \tol ) = 1\right]$
\end{enumerate}
\end{comment} 
\ITRconsistency

\newcommand{\PRELIMproofRCIF}{
\if1\prelim{
%rcif_consistency proof
\begin{comment}
\documentclass[12pt]{article}
\input{header} % Include the header file
\input{header2}
%\begin{document}
\textbf{Single Stage Random Cumulative Incidence Forest}
\end{comment}

{ %temporary command space
\newcommand{\shat}{\survnhat\argtleft}
\newcommand{\strue}{\truehsurv\argtleft}
\newcommand{\kdnj}{\kernel\argH d\Nj\argt}
\newcommand{\knj}{\kernel\argH \Nj\argt}
\newcommand{\ky}{\kernel\argH \atrisk \argt}
\newcommand{\tmpone}{\strue
\frac{\PnP \kdnj}{\Pn \ky}}

%("next" part)
\newcommand{\rn}{r_n}
%\newcommand{\sqrtn}{\sqrt{n}}
\newcommand{\thetanHAT}{\widehat{\theta}_n}
\newcommand{\truetheta}{\theta_0}
\newcommand{\fthat}{f(\thetanHAT)}
\newcommand{\fttrue}{f(\truetheta)}
\newcommand{\thetadiff}{(\thetanHAT - \truetheta)}
\newcommand{\fdot}{\dot{f}_{\truetheta}}
\newcommand{\fdotlong}{\dot{f}_{\truetheta + t(\thetanHAT-\truetheta)}}
\begin{proof}[Proof of Theorem \ifnum\jasa=1 3\else \ref{rcif_consistency}\fi]

\newcommand{\fracPn}{\frac{\PnkH d\Nj\argu}{\PnkH \atrisk\argu}}
\newcommand{\frachPn}{\frac{\PnkH d\Nj^\histlower\argu}{\PnkH \atrisk\argu}}

Our estimator is 
\begin{align*}
\AJhat\argtgh=\intzerotot \survnhhat\arguleft \frachPn.
\end{align*}

For simplicity, we drop the $\histlower$ superscript for any survival function when it does not cause confusion (i.e., $\survnhhat \equiv \Snhat$), but keep in mind that the survival functions in this proof use both arguments $t$ and $\histlower$, since the RCIF is built based on $\{\hist_{i}, \delta_{i}, \delta_{ji}\}_{i=1}^n$. We do the same for $\Nj\argt$ such that $\Nj^\histlower\argt \equiv \Nj\argt$ where $\Nj^\histlower\argt$ are the events due to cause-$j$ at time $t$ with history $\histlower$. 
\begin{align}
\label{eq:rcif_est}
\implies&\AJhat\argtgh\equiv\intzerotot \Snhat\arguleft \fracPn \notag \eqnspace
&=\intzerotot \Snhat\arguleft \fracPn 
- \intzerotot \truesurv\arguleft \fracPn \notag
\eqnspace
&\quad \quad + \intzerotot \truesurv\arguleft \fracPn \notag 
\eqnspace
&=\intzerotot \bigg\{ \Snhat\arguleft -\truesurv\arguleft \bigg\} \fracPn
+ \intzerotot \truesurv\arguleft \fracPn.
\end{align}
Let $\dLamnhHat\argt = \frac{\dLamnHATnum}{\dLamnHATdenom}$ and
$\dLamnhTilde\argt = \frac{\dLamnTILDEnum}{\dLamnTILDEdenom}$. First, we obtain a bound for $\bigg| \intzerotot \bigg\{ \Snhat\arguleft -\truesurv\arguleft \bigg\} \dLamnhHat\argt \bigg| $ for $t \in (0,\tau]$ and $\histlower \in \history$:
\begin{align*}
\left| \intzerotot \bigg\{ \survnhat\arguleft - \intzerotot \truesurv\arguleft \bigg\} d\LamnhHat\argu \right| 
\leq \bigg\{ \suptzerototauh
\left| \survnhat\argt -\truesurv\argt \right|  \bigg\}\intzerotot d\LamnhHat\argu,
%\arguleft is incorporated by \supzerototau
\end{align*}
where $\suptzerototau
\left| \survnhat\argt -\truesurv\argt \right|$ (with $t^{-}$ incorporated into $\suptzerototau$) has the rate for Kaplan-Meier for the single-stage generalized random survival forest, and $\intzerotot d\LamnhHat\args \leq \LamnhHat(\tau)$ is uniformly bounded over $\histlower$. Then, the first term of \cref{eq:rcif_est} converges to $0$.

\newcommand{\covrcif}{I\{s \ge t\} [N_0(s) - N_0(t^-)] - N_0(s)Y_0(t)}
\newcommand{\fracPnDivided}{\frac{\PnkH d\Nj\argu/\PnkH}{\PnkH \atrisk\argu/\PnkH}}

We next examine the second term of \cref{eq:rcif_est}. We can re-write it by dividing both the numerator and denominator by $\PnkH$ to get $\intzerotot \truesurv\arguleft \fracPnDivided$. We will use Lemma 12.3 from \cite{kosorok_introduction_2008}. %first verify $\{\delta_j I(X\leq t), t\ge0\}$ and $\{I(X\ge t), t\ge0\}$ are both Donsker, so that each component is Donsker. 
Since $\{\delta_j I(X\leq t), t\ge0\}$ and $\{I(X\ge t), t\ge0\}$ are indicators and $dN_j$ is at most $1$, they are bounded Donsker. Since $\Nj$ and $\atrisk$ are bounded Donsker, $\kernel$ is bounded Donsker, and the product of bounded Donsker is Donsker, then
\begin{align}
\label{eq:rcif_distr_0}
\sqrtn \, \bigmatlabel\arght \equiv 
    \bigmatlong
\rightsquigarrow
\begin{pmatrix}
    \Gone\arght\\
    \Gtwo\arght\\
    \Gthree\arght
\end{pmatrix},
\end{align}
where $\Gone$, $\Gtwo$, and $\Gthree$ are tight Gaussian processes. 
%We do this because we assume $\atrisk\argt$ is bounded away from $0$ but we need to show this is the case for $\kernel\argH$ in order to avoid $\Ah\argt \to 0$. 
Thus, $\sqrtn \, \bigmatlabel\arght 
= \Opht$, %uniform over \histlower and t
where $\Opht$ is a quantity $R_n\arght$ such that $\suptzerototauh \abs{R_n\arght} = \Op(1)$, which implies $\bigmatlabel\arght = \Ophtsqrtn$.

\begin{comment}
$$\sqrtn 
\begin{pmatrix}
\Pn \kernel\argH \Nj\argt - \PkH \Nj\argt\\
\Pn \kernel\argH \atrisk\argt - \PkH \atrisk \argt
\end{pmatrix} 
\rightsquigarrow
\begin{pmatrix}
    \Gone\arght\\
    \Gtwo\arght
\end{pmatrix},$$
\end{comment}

\newcommand{\invtruesurv}{\frac{1}{\truesurv\argtleft}}

Since $\truesurv\argtleft$ is bounded away from $0$, we have from \cref{eq:rcif_distr_0},
\begin{align}
\label{eq:rcif_distr}
    \sqrtn 
\begin{pmatrix}
\Pn \kernel\argH \Nj\argt - \PkH \Nj\argt\\
 \invtruesurv \Pn \kernel\argH \atrisk\argt - \invtruesurv \PkH \atrisk \argt\\
 \PnkH - \PkH
\end{pmatrix} 
\rightsquigarrow
\begin{pmatrix}
    \Gone\arght\\
   \invtruesurv \Gtwo\arght \\
   \Gthree\arght
\end{pmatrix}.
\end{align}

\begin{comment}
\documentclass[12pt]{article}
\input{header}
\input{header2}
%\begin{document}

Christina Zhou \& Michael Kosorok \\
Project 1 Theory
\section{Updates to RCIF Consistency Proof}
\end{comment}

{ %begin tmp space

%Remember, $\Ah\argt = \frac{\PnkH\atrisk\argt}{\truesurv\argtleft}$ and $\Bh\argt = \PnkH\Nj\argt$ where $\truesurv\argtleft$ is a constant function of $\histlower$ and $t$, but we drop the superscript $\histlower$ when notation is not confusing.
\begin{comment}
\begin{align*}
    \intzerotot \frac{d\Bh\argt}{\Ah\argt} - \intzerotot \frac{d\Bh_2\argt}{\Ah_2\argt}
\end{align*}
\end{comment}

First, we examine $\Nj$. From \cref{eq:rcif_distr_0}, we have
\begin{equation}
\label{rcif_nj_0}
    \frac{\PnkH\Nj\argt}{\PnkH} = \frac{\PkH\Nj\argt + \Ophtsqrtn}{\PkH + \Ophtsqrtn}.
\end{equation}
\newcommand{\tmprate}{n^{\beta-1}}
\newcommand{\tmprateinv}{n^{1-\beta}}
\newcommand{\tmpratetwo}{n^{\frac{1}{2}-\beta}}
\newcommand{\knot}{k_0}
\newcommand{\Ophtsqrtnbeta}{\Op^{\arght}(\tmpratetwo)}
\newcommand{\ophtsqrtnbeta}{\op^{\arght}(\tmpratetwo)}
\newcommand{\numun}{U_n}
%We now use the size of $\PkH$, so that this vanishes at the rate $\tmprate$.
Assumption \ifnum\jasa=1 5 \else \ref{assump:Cho1} \fi implies that $\inf_{\histlower \in \history} \PkH \ge \knot \tmprate$ for all $\histlower \in \history$ where $\knot > 0 $ is some constant. Note that $\tmprateinv \times n^{-1/2} = \tmpratetwo$. Multiplying $\tmprateinv$ to the numerator and denominator of the right-hand side of \cref{rcif_nj_0}, we obtain
\newcommand{\denomlong}{\tmprateinv \PkH + \Ophtsqrtnbeta}
\begin{align}
\frac{\PnkH\Nj\argt}{\PnkH} = &\frac{\tmprateinv \PkH\Nj\argt + \Ophtsqrtnbeta}{\denomlong}. \label{rcif_nj_1.1}
\end{align}
The denominator of the right-hand side of \cref{rcif_nj_1.1} is greater than or equal to $\knot + \opht$ because $\tmprateinv \PkH \ge \tmprateinv\knot\tmprate =\knot$ and the residual error is going to zero since by Assumption \ifnum\jasa=1 5\else \ref{assump:Cho1}\fi, $\beta>\frac{1}{2}$. Therefore, we have
\begin{equation}
    \begin{split}
\frac{\PnkH\Nj\argt}{\PnkH} = \frac{\tmprateinv\PkH\Nj\argt+\opht}{\tmprateinv\PkH+\opht}. \label{rcif_nj_1.2}
    \end{split}
\end{equation}
%Let $\numun = \tmprateinv \PkH\Nj\argt + \Ophtsqrtnbeta$.

We note that
\begin{equation}
\label{rcif_nj_2}
%\implies \frac{\numun}{\denomlong} &= \frac{\numun}{\tmprateinv \PkH} \bigg( 1 + \opht \bigg) \eqnspace
\frac{\tmprateinv \PkH + \opht}{\tmprateinv\PkH}
= 1+\frac{\opht}{\tmprateinv\PkH}
=1+\opht,
\end{equation}
where the last equality in \cref{rcif_nj_2} holds because $\tmprateinv\PkH \ge \knot > 0$.
Since $(1+\opht)^{-1}=1+\opht$, we can take the inverse of \cref{rcif_nj_2}.
\begin{align}
\implies \frac{\tmprateinv\PkH}{\tmprateinv \PkH + \opht} &= 1+\opht \notag \eqnspace
%divide each side by \tmprateinv\PkH
\implies \frac{1}{\tmprateinv \PkH + \opht} &= \frac{1 + \opht}{\tmprateinv \PkH}. \label{rcif_3}
\end{align}

To continue from \cref{rcif_nj_1.2}, multiply both sides of \cref{rcif_3} by $\tmprateinv \PkH\Nj\argt + \opht$.
\begin{align*}
%multiple each side by \tmprateinv \PkH\Nj\argt + \opht
&\implies \frac{\tmprateinv \PkH\Nj\argt + \opht}{\tmprateinv \PkH + \opht} = \frac{\tmprateinv \PkH\Nj\argt + \opht}{\tmprateinv \PkH} \bigg( 1 + \opht \bigg). %\notag \eqnspace
\end{align*}

We re-write the numerator as $\tmprateinv\PkH\Nj\argt=\tmpratetwo n^{\frac{1}{2}}\PkH\Nj\argt$. Since \\ $\sqrtn\PkH\Nj\argt = \Opht$ and $\tmpratetwo\to 0$, $\tmprateinv\PkH\Nj\argt$ is bounded. The denominator is bounded away from $0$ because $\tmprateinv\PkH \ge \knot>0$. This implies,
\begin{align*}
\frac{\PnkH\Nj\argt}{\PnkH} &=\frac{\tmprateinv \PkH\Nj\argt}{\tmprateinv \PkH} \bigg( 1 + \opht \bigg) + \opht %\label{residual}
\eqnspace
%can cancel out \tmprateinv now.
&=\frac{\PkH\Nj\argt}{\PkH} \bigg( 1 + \opht \bigg) + \opht \notag \eqnspace
%\frac{\PkH\Nj\argt}{\PkH} is bounded quantity, and little o_p times bounded quantity is still little o_p
&=\frac{\PkH\Nj\argt}{\PkH} + \opht \quad \text{(because $\frac{\PkH\Nj\argt}{\PkH}$ is bounded%, so $\frac{\PkH\Nj\argt}{\PkH} \times \opht = \opht$
).} 
%\label{lastline}
\end{align*}

Hence, $$\frac{\PnkH\Nj\argt}{\PnkH} = \frac{\PkH\Nj\argt}{\PkH} + \opht.$$ %and thus we have consistency of the kernel, or convergence for $\Nj\argt$.

Using similar arguments, we can show that the kernel version of $\atrisk$,  $\frac{1}{\truesurv\argtleft}\PnkH\atrisk\argt$ is consistent as well. Specifically, since $\truesurv\argt$ is bounded away from $0$, similar to \cref{rcif_nj_1.1}, we have
\begin{align*}
\frac{\frac{1}{\truesurv\argtleft}\PnkH\atrisk\argt}{\PnkH} = &\frac{\tmprateinv \frac{1}{\truesurv\argtleft}\PkH\atrisk\argt + \Ophtsqrtnbeta}{\denomlong}%\label{rcif_y_1.1} \eqnspace
&=\frac{\tmprateinv\frac{1}{\truesurv\argtleft}\PkH\atrisk\argt+\opht}{\tmprateinv\PkH+\opht}. %\label{rcif_y_1.2}
\end{align*}
Multiplying $\tmprateinv \frac{1}{\truesurv\argtleft}\PkH\atrisk\argt + \opht$ on the both sides of \cref{rcif_3} and simplifying the expressions, we have
\begin{align*}
%multiple each side by \tmprateinv \frac{1}{\truesurv\argtleft}\PkH\atrisk\argt + \opht
\implies \frac{\tmprateinv \frac{1}{\truesurv\argtleft}\PkH\atrisk\argt + \opht}{\tmprateinv \PkH + \opht} %&= \frac{\tmprateinv \frac{1}{\truesurv\argtleft}\PkH\atrisk\argt + \opht}{\tmprateinv \PkH} \bigg( 1 + \opht \bigg) \eqnspace
%&=\frac{\tmprateinv \frac{1}{\truesurv\argtleft}\PkH\atrisk\argt}{\tmprateinv \PkH} \bigg( 1 + \opht \bigg) + \opht \eqnspace
%%can cancel out \tmprateinv now.
%&=\frac{\frac{1}{\truesurv\argtleft}\PkH\atrisk\argt}{\PkH} \bigg( 1 + \opht \bigg) + \opht \eqnspace
%\frac{\frac{1}{\truesurv\argtleft}\PkH\atrisk\argt}{\PkH} is bounded quantity, and little o_p times bounded quantity is still little o_p
&=\frac{\frac{1}{\truesurv\argtleft}\PkH\atrisk\argt}{\PkH} + \opht.
\end{align*}
By Assumption \ifnum\jasa=1 5\else \ref{assump:CuiAssump1}\fi, $\PkH\atrisk\argt \ge \tauUB \PkH$. Then $\frac{\frac{1}{\truesurv\argtleft}\PkH\atrisk\argt}{\PkH} \ge \frac{\frac{1}{\truesurv\argtleft}\tauUB \PkH}{\PkH} = \frac{\tauUB}{\truesurv\argtleft}$. Since $\truesurv\argtleft$ is bounded away from $0$, we have $\frac{\frac{1}{\truesurv\argtleft}\PnkH\atrisk\argt}{\PnkH}$ is bounded away from $0$.
} % end of tmp space

\begin{comment}
% Bibliography
\bibliographystyle{apalike}
\bibliography{00_references/Overview, 00_references/00_references_manual}
%\end{document}
\end{comment}

\begin{comment}
%General Idea
Then, we can use Lemma 12.3 with continuous mapping theorem. Since $\Ah\argt$ is bounded below by $0$ and the map $\frac{1}{\Ah\argt}$ is smooth, we get $\frac{\truesurv\argtleft}{\atrisk
\argt}$. Because we have Hadamard differentiable functions, the map is continuous. We then want to use continuous mapping theorem, since each term is consistent so the continuous map is also consistent. With our assumptions, the continuous function of $\Bh\argt$ and $\Ah\argt$ gives the consistency of the random cumulative incidence forest.
\end{comment}

Let $\Ah\argu = \frac{1}{\truesurv\arguleft}\frac{\PnkH\atrisk\argu}{\PnkH}$ and $\Bh\argu = \frac{\PnkH\Nj\argu}{\PnkH}$. Let $\Bh_0\argt= \frac{\PkH \Nj\argt}{\PkH}$ and $\Ah_0\argt = \invtruesurv \frac{\PkH \atrisk \argt}{\PkH}$. Then our point of interest is $(B,A) = (\Bh_0\argt,\Ah_0\argt)$. The components of the second term of our estimator $\AJhat\argtgh$ in \cref{eq:rcif_est} can be mapped twice:
\begin{align*}
    (B,A) \mapsto (B,\frac{1}{A}) \mapsto \int_{[0,t]} \frac{1}{\Ah\argu}\,d\Bh \argu.
\end{align*}

Using Lemma 12.3 and Lemma 6.19 from \cite{kosorok_introduction_2008}, this composition map is Hadamard differentiable on a domain of type 
$\{ (B,A) :
\suph\int_{[0,\tau]} \lvert d\Bh\argt \rvert \leq M,\\
\inf_{\substack{t \in [0, \tau] \\ \histlower \in \history}} \lvert \Ah\argt \rvert \geq \epsilon \}$ 
for all $\histlower \in \history$ and for a given \( M < \infty \) and \( \epsilon > 0 \), at every point \((B,A)\) such that \( \frac{1}{\Ah\argt} \) has bounded variation. Thus, the map is continuous. Using Assumption \ifnum\jasa=1 9\else \ref{khbddtv}\fi, the total variation of the kernels is bounded. Since we have shown that $\Ah\argt$ and $\Bh\argt$ are consistent, and $\Ah\argt$ is bounded below by $0$ and the map $\frac{1}{\Ah\argt}$ is smooth, we apply the continuous mapping theorem so that the continuous map is also consistent.

Thus,
\begin{equation*}
\suptzerototauh \abs{\AJhat \argtgh - \cuminc_j \argtgh} \to 0
\end{equation*}
in probability as $n \to \infty$.

\begin{comment}
Then by Lemma 12.3, the derivative of the composition map is given by $$ (\alpha,\beta) \mapsto (\alpha, \frac{-\beta\truesurv^2}{\atrisk_0^2}) \mapsto \int_{[0,t]} \frac{\truesurv d\alpha}{\atrisk_0} - \int_{[0,t]} \frac{\beta \truesurv dN_0}{\atrisk_0^2}$$.

Then, from \cref{eq:rcif_distr} and using Theorem 2.8 in \cite{kosorok_introduction_2008}, we obtain \begin{align*}
    \sqrtn (\AJhat - \cuminc_j) \rightsquigarrow \int_{[0,(\cdot)]} \frac{\truesurv d\Gone}{\atrisk_0} - \int_{[0,(\cdot)]} \frac{\truesurv\Gtwo dN_0}{\atrisk_0^2}
\end{align*}
\end{comment}
\begin{comment}
Next, we obtain a bound for $\dLamnhHat$ vs. $\dLamnhTilde$ for $t \in (0,\tau]$. Using Riemann-Stieltjes integration by parts, we have
\begin{align*}
\intzerotot \truesurv\argsleft &\left[ \dLamnhHat\args - \dLamnhTilde\args \right] \eqnspace
&= \truesurv\argtleft \left[\LamnhHat\argt - \LamnhTilde\argt\right] - \truesurv(0) \left[ \LamnhHat(0) - \LamnhTilde(0)
\right] 
-\intzerotot \left[ \LamnhHat\args - \LamnhTilde\args \right] d\truesurv\args
\eqnspace
&= \truesurv\argtleft \left[\LamnhHat\argt - \LamnhTilde\argt\right] 
-\intzerotot \left[ \LamnhHat\args - \LamnhTilde\args \right] d\truesurv\args
\eqnspace
&\leq 2 \suptzerototauh \left| \LamnhHat\argt - \LamnhTilde\argt \right|
%this has the rate from the derivation from Lemma 12.3?
\end{align*}

For Lemma 12.3 and bounded kernels (also verifying that total variance is bounded), we have:
\begin{align*}
    \sqrt{n}
\begin{pmatrix}
\survnhhat\argt - \truehsurv\argt \\
\Pn \kernel\argH \Nj\argt - \PkH \Nj\argt\\
\Pn \kernel\argH \atrisk\argt - \PkH \atrisk \argt
\end{pmatrix}
= \sqrt{n}
\begin{pmatrix}
\survnhhat\argt - \truehsurv\argt \\
\PnP \knj\\
\PnP \ky
\end{pmatrix}
= \Opht %uniform over \histlower and t
\end{align*}

 Then,
\begin{align}
&\intzerotou \shat \dLamnHAT - \intzerotou \strue \dLamnTILDE = \intzerotou \shat\frac{\dLamnHATnum}{\dLamnHATdenom} - \intzerotou \strue \frac{\dLamnTILDEnum}{\dLamnTILDEdenom} \notag \eqnspace
&= \intzerotou \shat \dLamnHAT - \strue \dLamnTILDE + \left[ \strue \dLamnHAT - \strue \dLamnHAT \right] \notag \eqnspace
&= \intzerotou \left[ \shat - \strue \right] \dLamnHAT \\
&\quad \quad \quad + \intzerotou \tmpone - \intzerotou \strue 
\left\{\frac{\PnP \ky}{\dLamnHATdenom}\right\}\dLamnTILDE, \notag
\end{align}
because $\strue\dLamnHAT - \strue\dLamnTILDE$
\begin{align*}
%\strue\dLamnHAT &- \strue\dLamnTILDE \eqnspace
&= \strue\left[\frac{\dLamnHATnum-\dLamnTILDEnum+\dLamnTILDEnum}{\dLamnHATdenom} - \frac{\dLamnTILDEnum}{\dLamnTILDEdenom}\right] \eqnspace 
&= \strue\left[ 
\frac{\PnP \kdnj}{\Pn \ky} + \frac{\dLamnTILDEnum}{\dLamnHATdenom} - \frac{\dLamnTILDEnum}{\dLamnTILDEdenom}
\right]\eqnspace
&= \strue\left[ 
\frac{\PnP \kdnj}{\Pn \ky} - \frac{\dLamnTILDEnum}{\dLamnTILDEdenom} + \frac{\dLamnTILDEnum}{\dLamnHATdenom} 
\right]\eqnspace
&= \tmpone - \strue \left\{\frac{\dLamnTILDEnum}{\dLamnTILDEdenom} - \frac{\dLamnTILDEnum}{\dLamnHATdenom} 
\right\}\eqnspace
&= \tmpone - \strue 
\dLamnTILDEnum
\left\{\frac{\dLamnHATnum - \dLamnTILDEnum}{\dLamnTILDEdenom\dLamnHATdenom} 
\right\}\eqnspace
&= \tmpone - \strue 
\left\{\frac{\PnP \ky}{\dLamnHATdenom}\right\}\frac{\dLamnTILDEnum}{\dLamnTILDEdenom} \eqnspace
&= \tmpone - \strue 
\left\{\frac{\PnP \ky}{\dLamnHATdenom}\right\}\dLamnTILDE
\end{align*}

%next part
We know $\sqrtn \thetadiff = \widetilde{O}_p(1)$. Let $\fthat - \fttrue$ be a Hadamard differentiable function. We note that $\thetanHAT$ could be infinite dimensional.

Using Lemma 12.3 from Kosorok 2008 Springer, we can show that
\begin{align*}
\sqrtn \left(\fthat - \fttrue\right) &= \fdot \left( \sqrtn \thetadiff \right) + \op \eqnspace
&= \int^1_0 \fdotlong \left(\sqrtn \thetadiff\right)dt \eqnspace
&=\fdot \left( \sqrtn \thetadiff \right) + \int^1_0 \left[ \fdotlong - \fdot \right]dt \sqrtn \thetadiff
\end{align*}

In other words, $\fdotlong - \fdot \to 0$ within the bounded set $\thetadiff$ when $0<t<1$. We have convergence as $\sqrtn \thetadiff = \widetilde{O}_p(1) \implies \rn \norm{\thetanHAT-\truetheta} = \Opht$.
\end{comment}

\end{proof}
} %end command space

\begin{comment}
% Bibliography
\bibliographystyle{apalike}
\bibliography{00_references/Overview, 00_references/00_references_manual}
\end{document}
\end{comment}

} \fi
}


\newcommand{\PRELIMproofITR}{
\if1\prelim{
%The proof for Theorem \ref{thm:ITR_Consistency} is below.
%%%%%%%%%%%%%%%%%%%%%%%%%%%%%%%%%%%%%%%%%%%%%%%%%%%%%%%%%%%
%%%%%%%%%%%%%%%%%%%%%%%%%%%%%%%%%%%%%%%%%%%%%%%%%%%%%%%%%%%
%%%%%%%%%%%%%%%%%%%%%%%%%%%%%%%%%%%%%%%%%%%%%%%%%%%%%%%%%%%
% Consistency of ITR estimator
\newcommand{\ITRconsistency}{
\begin{theorem} (Consistency of the individualized treatment regime estimator) 
\label{thm:ITR_Consistency}
Assume Assumptions \ref{assump:CuiAssump1}-\ref{assump:CuiAssump2} hold. Let $\doptHAT$ denote the optimal ITR estimator that is built following Algorithm \ref{alg:Overview} and Assumptions \ref{assump:boundeddensity}-\ref{assump:boundedV0}. Assume Assumptions \ref{assump:Causal1}-\ref{assump:Causal3} hold and that the random forests are built for each treatment arm. Then, for %each $\phi = \phi^{\mu}, \phi^{\sigma,t_0}$ where 
$\phi^\mu = \begin{pmatrix}
    \phi^\mu_1 \\ \phi^\mu_2
\end{pmatrix},$ % and $\phi^{\sigma, t_0} = \begin{pmatrix} \phi^{\sigma, t_0}_1 \\ \phi^{\sigma, t_0}_2 \end{pmatrix},$
\[ \lim_{n \to \infty} \| \mathbb{V}(\doptTRUE(Z)) - \mathbb{V}(\doptHAT(Z)) \| \xrightarrow{p} 0\]
where $\|\cdot\|$ is defined as the $L_2$ norm.
\end{theorem}
}

\newcommand{\ITRconsistencyPROOF}{
\begin{proof}[Proof of Theorem \ifnum\jasa=1 1\else \ref{thm:ITR_Consistency}\fi]
%We know that, by properties of norm, $\left\| \begin{pmatrix} a_1 \\ a_2 \end{pmatrix} + \begin{pmatrix} b_1 \\ b_2 \end{pmatrix}\right\| \leq \left\| \begin{pmatrix} a_1 \\ a_2 \end{pmatrix} \right\| + \left\| \begin{pmatrix} b_1 \\ b_2 \end{pmatrix} \right\|$ and $\left\| c\begin{pmatrix} a_1 \\ a_2 \end{pmatrix} \right\| =  |c| \left\|\begin{pmatrix} a_1 \\ a_2 \end{pmatrix}\right\|$ where $c$ is a scalar. Also, by property of expectation, we know that $\|\mathbb{E}(\cdot)\| \leq \mathbb{E}(\|\cdot\|)$.

Note that
\begin{equation*}
\begin{split}
\left\| \mathbb{V}(\doptTRUE(Z)) - \mathbb{V}(\doptHAT(Z)) \right\| 
% line 1
&= \left\|\sum_{a \in \mathcal{A}} \mathbb{E}_Z\left[ \valueTRUE
\cdot \left( \mathbf{1}\{\doptTRUE(Z) = a\} - \mathbf{1}\{\doptHAT(Z) = a\} \right)
\right] \right\| \\
% line 2
&\leq \sum_{a \in \mathcal{A}} \left\|\mathbb{E}_Z\left[ \valueTRUE
\cdot \left( \mathbf{1}\{\doptTRUE(Z) = a\} - \mathbf{1}\{\doptHAT(Z) = a\} \right)
\right] \right\| \\
% line 3
&\leq \sum_{a \in \mathcal{A}} \mathbb{E}_Z \left\| \left[ \valueTRUE
\cdot \left( \mathbf{1}\{\doptTRUE(Z) = a\} - \mathbf{1}\{\doptHAT(Z) = a\} \right)
\right] \right\| \\
% line 4
&\leq \sum_{a \in \mathcal{A}} \mathbb{E}_Z \left[ \left\| \valueTRUE \right\|
\cdot \left| \mathbf{1}\{\doptTRUE(Z) = a\} - \mathbf{1}\{\doptHAT(Z) = a\} \right|
\right] \\
% line 5
&= \mathbb{E}_Z \left[ \sum_{a \in \mathcal{A}} \left( \mathbf{1} \{\doptTRUE(Z) \neq \doptHAT(Z)\} \left\|\valueTRUE\right\| \right) \right] \\
% line 6
&\leq \mathbb{E}_Z \left( \mathbf{1} \{\doptTRUE(Z) \neq \doptHAT(Z)\} \right) \sup_{a,Z} \left\| \valueTRUE \right\| \\
% line 7
&= \text{Pr}(\doptTRUE(Z) \neq \doptHAT(Z) | \doptHAT) \sup_{a,Z} \left\| \valueTRUE \right\|\\
% line 8
&\leq \sevenbound \text{Pr}(\doptTRUE(Z) \neq \doptHAT(Z) | \doptHAT).\\
\end{split} % split gives one number; align gives a number per line
\end{equation*} The third to last line holds because 1) only one term in the summation will have a non-zero contribution and 2) the supremum represents the ``worst case scenario" over all fixed combinations of $a$ and $z$ and thus does not depend on the random variables. The last line holds because we assumed $\sup_{a,Z} \left\| \valueTRUE \right\|$ is bounded by $\sevenbound$ via Assumption \ifnum\jasa=1 11\else \ref{assump:boundedV0} \fi. Thus, the norm of this distance, $\mathbb{V}(\doptTRUE) - \mathbb{V}(\doptHAT)$, goes to $0$ in probability if $\text{Pr}(\doptTRUE(Z) \neq \doptHAT(Z)|\hat{d}_\text{opt}) \xrightarrow{p} 0$.\\

We complete the proof by first showing
\begin{align}
    \label{eqn:step1}
        \sum_{a \in \mathcal{A}} \left| I\{ \valueoneTRUE - (1-\tol)\valueonemaxTRUE \ge 0 \} - I\{ \valueoneHAT - (1-\tol)\valueonemaxHAT \ge 0 \} \right| \rightarrow 0.
        \end{align}
    %(Note: we take the difference because both are random (one relies on Z and the other is an estimator)).
The proof of \cref{eqn:step1} will be presented shortly. Now, \cref{eqn:step1} implies 
\begin{align}
    \label{eqn:star}
        \argmax_{a \in \mathcal{A}: \valueoneHAT \ge (1-\tol)\valueonemaxHAT} \valuetwoHAT \xrightarrow{p} \argmax_{a \in \mathcal{A}: \valueoneTRUE \ge (1-\tol)\valueonemaxTRUE} \valuetwoHAT.
    \end{align} 
Theorem \ifnum \jasa=1 3 \else \ref{thm:RCIF_Consistency} \fi in our paper states that RCIFs are consistent. %replace <- with theorem when typed up.
Thus, by the Continuous Mapping Theorem, $\valuetwoHAT \xrightarrow{p} \valuetwoTRUE$. %$\valuetwoHAT$ is bounded.
Then,
\begin{align}
\label{eqn:step2}
        \sup_{a \in \mathcal{A}} \, \mathbb{E} \, \left| \valuetwoHAT - \valuetwoTRUE \right| \xrightarrow{p} 0,
    \end{align} 
which implies that the right hand side of \cref{eqn:star} converges to \\$\argmax_{a \in \mathcal{A}: \valueoneTRUE \ge (1-\tol)\valueonemaxTRUE} \valuetwoTRUE$. Thus, since $Z$ is random, $$\text{Pr}(\doptTRUE(Z) \neq \doptHAT(Z) | \doptHAT) \xrightarrow{p} 0.$$ This completes the proof of Theorem \ifnum\jasa=1 1\else \ref{thm:ITR_Consistency}\fi.
\end{proof}
}

\newcommand{\eqnsteponePROOF}{
%\textbf{Proof of \cref{eqn:step1}:}
\begin{proof}[Proof of \cref{eqn:step1}]
Let $A_0 = \valueoneTRUE - (1-\tol)\valueonemaxTRUE$
and $\widehat{A}_n = \valueoneHAT - (1-\tol)\valueonemaxHAT$. These are functions of random variable $Z$. Then, we can re-write \cref{eqn:step1} as 
$$\sum_{a \in \mathcal{A}} \left| I\{ A_0 \ge 0 \} - I \{ \widehat{A}_n \ge 0 \} \right| \rightarrow 0.$$
\begin{equation}
\label{eqn:terms12}
\begin{split}
    \left| I\{ \widehat{A}_n \ge 0 \} - I \{A_0 \ge 0 \} \right| 
    &= I\{ \widehat{A}_n < 0 \le A_0 \} + I \{A_0 < 0 \le \widehat{A}_n \} \\
    &= I\{ \widehat{A}_n - A_0 < -A_0 \leq 0 \} + I \{0 < -A_0 \le \widehat{A}_n - A_0 \} \\
    &= I\{ 0 \le A_0 < -(\widehat{A}_n - A_0) \} + I \{- (\widehat{A}_n - A_0) \le A_0 < 0 \}. \\
\end{split}
\end{equation}
In \cref{eqn:terms12}, the first equality holds because these are the only possible indicators such that their sum is equal to $1$. We get the 2nd equality by subtracting $A_0$ on both sides within the indicator functions. The third equality is due to multiplying by $-1$ and flipping the inequalities within the indicators.

Since we know $\valueoneHAT \xrightarrow{p} \valueoneTRUE$ and $\valueonemaxHAT \xrightarrow{p} \valueonemaxTRUE$ as shown in \cite{cho_multi-stage_2023}; and $\widehat{A}_n$ and $A_0$ are linear combinations of $\valueoneHAT$ and $\valueonemaxHAT$, and $\valueoneTRUE$ and $\valueonemaxTRUE$, respectively, then %Theorem \ref{thm:consistency_lincomb},
$\widehat{A}_n \xrightarrow{p} A_0$ because the linear combination of consistent estimators is consistent. This means
\begin{align}
\label{eqn:AnA0} 
\widehat{A}_n - A_0 \xrightarrow{p} 0.
\end{align}

By definition of convergence in probability, \cref{eqn:AnA0} implies $\exists\epsilon_n: Pr(|\widehat{A}_n - A_0|>\epsilon_n) \rightarrow 0 \text{ as } \epsilon_n \rightarrow 0$ for a fixed sequence $\epsilon_n$. Then, using Assumption \ifnum\jasa=1 10 \else \ref{assump:boundeddensity} \fi and \cref{eqn:AnA0}, we can bound \cref{eqn:terms12}. Starting with the first term of \cref{eqn:terms12},
\begin{equation*}
\label{eqn:term1}
\begin{split}
&I\{ 0 \le A_0 < -(\widehat{A}_n - A_0), |\widehat{A}_n-A_0|\le \epsilon_n \}  + I\{ 0 \le A_0 < -(\widehat{A}_n - A_0), |\widehat{A}_n-A_0| > \epsilon_n \} \le \\&I\{ 0 \le A_0 < \epsilon_n\} + I \{ |\widehat{A}_n - A_0| > \epsilon_n \}, \\    
\end{split}
\end{equation*}
 where $I\{ 0 \le A_0 < \epsilon_n\} \rightarrow 0$ by Assumption \ifnum\jasa=1 10 \else \ref{assump:boundeddensity}\fi (bounded density assumption) and
$I \{ |\widehat{A}_n - A_0| > \epsilon_n \} \rightarrow 0$ by \cref{eqn:AnA0} (and definition of consistency). Specifically, assuming Assumption \ifnum\jasa=1 10\else \ref{assump:boundeddensity}\fi,
$Pr(0\le A_0 \le \epsilon_n) = \int^{\epsilon_n}_{0} f_{A_0}(a_0)da_0 \approx f_{a_0}(0)\epsilon_n \le b\epsilon_n \rightarrow 0 $ as $\epsilon_n \rightarrow 0 $ for some bound $b$. The second term can be shown to be bounded similarly. 
\begin{equation*}
\begin{split}
&I \{- (\widehat{A}_n - A_0) \le A_0 < 0 \} = \\
&I \{- (\widehat{A}_n - A_0) \le A_0 < 0 , |\widehat{A}_n - A_0| \le \epsilon_n\} +
I \{- (\widehat{A}_n - A_0) \le A_0 < 0 , |\widehat{A}_n - A_0| > \epsilon_n\} \leq\\
&I(-\epsilon_n \leq A_0 < 0) + I(|\widehat{A}_n - A_0| > \epsilon_n).
\end{split}
\end{equation*} The first term for the inequality is because $-(\widehat{A}_n - A_0) < 0$, $\widehat{A}_n - A_0$ must be positive; thus, $-\epsilon_n \leq \widehat{A}_n - A_0 \leq \epsilon_n$.  The second term is because the indicator is a subset.

Because of Equation \cref{eqn:AnA0}, $Pr(|\widehat{A}_n - A_0| > \epsilon_n) \rightarrow 0$, and $Pr(-\epsilon_n \leq A_0 < 0) \rightarrow 0$ due to Assumption \ifnum\jasa=1 10\else \ref{assump:boundeddensity}\fi.
Specifically, $I(-\epsilon_n \leq A_0 < 0)  \implies Pr(-\epsilon_n \le A_0 < 0) \approx \epsilon_n f_{a_0}(0) < \epsilon_n b \rightarrow 0$ as $\epsilon_n \rightarrow 0$.
This completes the proof of \cref{eqn:step1}.
\end{proof}
}

\begin{comment}
MARCH 2024 UPDATE: WE CAN IGNORE BELOW I THINK.

We want to use the below equation. I think he said that, as long as the density of the quantity $\valueoneTRUE - (1-\tol)\valueonemaxTRUE$ is bounded then the whole thing converges to 0?
\begin{equation}
%\begin{align}
    \lim_{n \to \infty} E\left(\left|I( \valueoneHAT - (1-\tol)\valueonemaxHAT \ge 0) - I( \valueoneTRUE - (1-\tol)\valueonemaxTRUE \ge 0 )\right|\right) = 0 \\
%\end{align}
\end{equation}
We want to show that $\valueoneHAT - (1-\tol)\valueonemaxHAT -\xrightarrow{p} \valueoneTRUE - (1-\tol)\valueonemaxTRUE$
Let $X_n \xrightarrow{p} \mu$ and $Y_n \xrightarrow{p} \gamma $. Then, $Pr(|X_n - \mu| > \epsilon \text{ or } |Y_n - \gamma|>\epsilon) \leq Pr(|X_n - \mu|>\epsilon) + Pr(|Y_n - \gamma|>\epsilon)$. If $n \rightarrow \infty$, then both $Pr(|X_n - \mu|>\epsilon) \rightarrow 0$ and $Pr(|Y_n - \gamma|>\epsilon) \rightarrow 0$. We want to use this as the generic idea for $X_n$ being OS and $Y_n$ being CIF.
\end{comment}

\begin{comment}
MARCH 2024 UPDATE: DEFINITELY IGNORE BELOW. PROOF IS COMPLETE.

\subsubsection{OLD STUFF: Notes from before 2/20/24 meeting}

Sample size $n$ such that $\hat{d}_n$. Let 
\[
V(\hat{d}_n) = \begin{pmatrix}
V_1(\hat{d}_n) \\ V_2(\hat{d}_n)
\end{pmatrix}
\]

We want to show:
\[
\lim_{{n \to \infty}} \bigg( \frac{V_1(d_*) - V_1(\hat{d}_n)}{V_1(d_*)} - \tol \bigg)_+ \xrightarrow{p} 0
\]
\\

Our optimal $d_{**}$ satisfies the following two criteria:
\begin{enumerate}
    \item $E\left[ I(\frac{V_1(d_*|Z)-V_1(d_{**}|Z)}{V_1(d_*|Z)} \leq \tol ) \right] = 1$
    \item maximize $E\left[V_2(d_{**}|Z) \;\middle|\; I(\frac{V_1(d_*|Z)-V_1(d_{**}|Z)}{V_1(d_*|Z)} \leq \tol ) = 1\right]$
\end{enumerate}
\end{comment}
\ITRconsistencyPROOF
\eqnsteponePROOF
} \fi
}


\PRELIMproofITR

\subsection{Random Forests}
\label{Sec:Forest_Consistency}
Consistency of our proposed estimator depends on consistency of the forests. We show this in Theorems \ref{genRSF_consistency} and \ref{rcif_consistency}. For simplicity, let $h = (a,Z)$.

% Theorem
\newcommand{\genrsfconsistency}{
\begin{theorem}[Uniform consistency of a random survival forest]
\label{genRSF_consistency}
 Suppose that Assumptions \ref{assump:CuiAssump1}-\ref{assump:CuiAssump2} hold. Let $\survfunchat\argtgh$ be an estimator of the survival probability $\survfunc\argtgh$ that is uniformly consistent in both $\histlower \in \history$ and $t \in [0, \tau]$. Then, the single-stage version of the random survival forest $\survfunchat\argtgh$ in \cite{cho_multi-stage_2023} built based on $\{\hist_{i}, \delta_{i}\}_{i=1}^n$ is uniformly consistent. Specifically,
\[
\suptzerototauh \abs{\survfunchat \argtgh - \survfunc \argtgh} \to 0
\]
in probability as $n \to \infty$.
\end{theorem}
}

\newcommand{\rcifconsistency}{
\begin{theorem}[Uniform consistency of a random cumulative incidence forest]
\label{rcif_consistency}
Suppose that bounded kernels for the random cumulative incidence forest (RCIF) hold and have total bounded variation (Assumption \ref{khbddtv}). Let $\AJhat\argtgh$ be an estimator of the cause-$j$ cumulative incidence function $\cuminc_j\argtgh$ that is uniformly consistent in both $\histlower \in \history$ and $t \in [0, \tau]$. Then, the RCIF $\AJhat(t \mid \histlower)$ built based on $\{\hist_{i}, \delta_{i}, \delta_{ji}\}_{i=1}^n$ is uniformly consistent. Specifically,
\[
\suptzerototauh \abs{\AJhat \argtgh - \cuminc_j \argtgh} \to 0
\]
in probability as $n \to \infty$.
\end{theorem}
}

\begin{comment}
\documentclass[12pt]{article}
\input{header} % Include the header file
\input{header2}
\begin{document}
\rcifconsistency
\genrsfconsistency
\end{document}
\end{comment}
\genrsfconsistency
\if1\prelim
{
\proofSM
} \fi

\rcifconsistency

\ifnum\prelim=1
\ifnum\jasa=0
% proof included for prelim
\newcommand{\PRELIMproofRCIF}{
\if1\prelim{
%rcif_consistency proof
\begin{comment}
\documentclass[12pt]{article}
\input{header} % Include the header file
\input{header2}
%\begin{document}
\textbf{Single Stage Random Cumulative Incidence Forest}
\end{comment}

{ %temporary command space
\newcommand{\shat}{\survnhat\argtleft}
\newcommand{\strue}{\truehsurv\argtleft}
\newcommand{\kdnj}{\kernel\argH d\Nj\argt}
\newcommand{\knj}{\kernel\argH \Nj\argt}
\newcommand{\ky}{\kernel\argH \atrisk \argt}
\newcommand{\tmpone}{\strue
\frac{\PnP \kdnj}{\Pn \ky}}

%("next" part)
\newcommand{\rn}{r_n}
%\newcommand{\sqrtn}{\sqrt{n}}
\newcommand{\thetanHAT}{\widehat{\theta}_n}
\newcommand{\truetheta}{\theta_0}
\newcommand{\fthat}{f(\thetanHAT)}
\newcommand{\fttrue}{f(\truetheta)}
\newcommand{\thetadiff}{(\thetanHAT - \truetheta)}
\newcommand{\fdot}{\dot{f}_{\truetheta}}
\newcommand{\fdotlong}{\dot{f}_{\truetheta + t(\thetanHAT-\truetheta)}}
\begin{proof}[Proof of Theorem \ifnum\jasa=1 3\else \ref{rcif_consistency}\fi]

\newcommand{\fracPn}{\frac{\PnkH d\Nj\argu}{\PnkH \atrisk\argu}}
\newcommand{\frachPn}{\frac{\PnkH d\Nj^\histlower\argu}{\PnkH \atrisk\argu}}

Our estimator is 
\begin{align*}
\AJhat\argtgh=\intzerotot \survnhhat\arguleft \frachPn.
\end{align*}

For simplicity, we drop the $\histlower$ superscript for any survival function when it does not cause confusion (i.e., $\survnhhat \equiv \Snhat$), but keep in mind that the survival functions in this proof use both arguments $t$ and $\histlower$, since the RCIF is built based on $\{\hist_{i}, \delta_{i}, \delta_{ji}\}_{i=1}^n$. We do the same for $\Nj\argt$ such that $\Nj^\histlower\argt \equiv \Nj\argt$ where $\Nj^\histlower\argt$ are the events due to cause-$j$ at time $t$ with history $\histlower$. 
\begin{align}
\label{eq:rcif_est}
\implies&\AJhat\argtgh\equiv\intzerotot \Snhat\arguleft \fracPn \notag \eqnspace
&=\intzerotot \Snhat\arguleft \fracPn 
- \intzerotot \truesurv\arguleft \fracPn \notag
\eqnspace
&\quad \quad + \intzerotot \truesurv\arguleft \fracPn \notag 
\eqnspace
&=\intzerotot \bigg\{ \Snhat\arguleft -\truesurv\arguleft \bigg\} \fracPn
+ \intzerotot \truesurv\arguleft \fracPn.
\end{align}
Let $\dLamnhHat\argt = \frac{\dLamnHATnum}{\dLamnHATdenom}$ and
$\dLamnhTilde\argt = \frac{\dLamnTILDEnum}{\dLamnTILDEdenom}$. First, we obtain a bound for $\bigg| \intzerotot \bigg\{ \Snhat\arguleft -\truesurv\arguleft \bigg\} \dLamnhHat\argt \bigg| $ for $t \in (0,\tau]$ and $\histlower \in \history$:
\begin{align*}
\left| \intzerotot \bigg\{ \survnhat\arguleft - \intzerotot \truesurv\arguleft \bigg\} d\LamnhHat\argu \right| 
\leq \bigg\{ \suptzerototauh
\left| \survnhat\argt -\truesurv\argt \right|  \bigg\}\intzerotot d\LamnhHat\argu,
%\arguleft is incorporated by \supzerototau
\end{align*}
where $\suptzerototau
\left| \survnhat\argt -\truesurv\argt \right|$ (with $t^{-}$ incorporated into $\suptzerototau$) has the rate for Kaplan-Meier for the single-stage generalized random survival forest, and $\intzerotot d\LamnhHat\args \leq \LamnhHat(\tau)$ is uniformly bounded over $\histlower$. Then, the first term of \cref{eq:rcif_est} converges to $0$.

\newcommand{\covrcif}{I\{s \ge t\} [N_0(s) - N_0(t^-)] - N_0(s)Y_0(t)}
\newcommand{\fracPnDivided}{\frac{\PnkH d\Nj\argu/\PnkH}{\PnkH \atrisk\argu/\PnkH}}

We next examine the second term of \cref{eq:rcif_est}. We can re-write it by dividing both the numerator and denominator by $\PnkH$ to get $\intzerotot \truesurv\arguleft \fracPnDivided$. We will use Lemma 12.3 from \cite{kosorok_introduction_2008}. %first verify $\{\delta_j I(X\leq t), t\ge0\}$ and $\{I(X\ge t), t\ge0\}$ are both Donsker, so that each component is Donsker. 
Since $\{\delta_j I(X\leq t), t\ge0\}$ and $\{I(X\ge t), t\ge0\}$ are indicators and $dN_j$ is at most $1$, they are bounded Donsker. Since $\Nj$ and $\atrisk$ are bounded Donsker, $\kernel$ is bounded Donsker, and the product of bounded Donsker is Donsker, then
\begin{align}
\label{eq:rcif_distr_0}
\sqrtn \, \bigmatlabel\arght \equiv 
    \bigmatlong
\rightsquigarrow
\begin{pmatrix}
    \Gone\arght\\
    \Gtwo\arght\\
    \Gthree\arght
\end{pmatrix},
\end{align}
where $\Gone$, $\Gtwo$, and $\Gthree$ are tight Gaussian processes. 
%We do this because we assume $\atrisk\argt$ is bounded away from $0$ but we need to show this is the case for $\kernel\argH$ in order to avoid $\Ah\argt \to 0$. 
Thus, $\sqrtn \, \bigmatlabel\arght 
= \Opht$, %uniform over \histlower and t
where $\Opht$ is a quantity $R_n\arght$ such that $\suptzerototauh \abs{R_n\arght} = \Op(1)$, which implies $\bigmatlabel\arght = \Ophtsqrtn$.

\begin{comment}
$$\sqrtn 
\begin{pmatrix}
\Pn \kernel\argH \Nj\argt - \PkH \Nj\argt\\
\Pn \kernel\argH \atrisk\argt - \PkH \atrisk \argt
\end{pmatrix} 
\rightsquigarrow
\begin{pmatrix}
    \Gone\arght\\
    \Gtwo\arght
\end{pmatrix},$$
\end{comment}

\newcommand{\invtruesurv}{\frac{1}{\truesurv\argtleft}}

Since $\truesurv\argtleft$ is bounded away from $0$, we have from \cref{eq:rcif_distr_0},
\begin{align}
\label{eq:rcif_distr}
    \sqrtn 
\begin{pmatrix}
\Pn \kernel\argH \Nj\argt - \PkH \Nj\argt\\
 \invtruesurv \Pn \kernel\argH \atrisk\argt - \invtruesurv \PkH \atrisk \argt\\
 \PnkH - \PkH
\end{pmatrix} 
\rightsquigarrow
\begin{pmatrix}
    \Gone\arght\\
   \invtruesurv \Gtwo\arght \\
   \Gthree\arght
\end{pmatrix}.
\end{align}

\begin{comment}
\documentclass[12pt]{article}
\input{header}
\input{header2}
%\begin{document}

Christina Zhou \& Michael Kosorok \\
Project 1 Theory
\section{Updates to RCIF Consistency Proof}
\end{comment}

{ %begin tmp space

%Remember, $\Ah\argt = \frac{\PnkH\atrisk\argt}{\truesurv\argtleft}$ and $\Bh\argt = \PnkH\Nj\argt$ where $\truesurv\argtleft$ is a constant function of $\histlower$ and $t$, but we drop the superscript $\histlower$ when notation is not confusing.
\begin{comment}
\begin{align*}
    \intzerotot \frac{d\Bh\argt}{\Ah\argt} - \intzerotot \frac{d\Bh_2\argt}{\Ah_2\argt}
\end{align*}
\end{comment}

First, we examine $\Nj$. From \cref{eq:rcif_distr_0}, we have
\begin{equation}
\label{rcif_nj_0}
    \frac{\PnkH\Nj\argt}{\PnkH} = \frac{\PkH\Nj\argt + \Ophtsqrtn}{\PkH + \Ophtsqrtn}.
\end{equation}
\newcommand{\tmprate}{n^{\beta-1}}
\newcommand{\tmprateinv}{n^{1-\beta}}
\newcommand{\tmpratetwo}{n^{\frac{1}{2}-\beta}}
\newcommand{\knot}{k_0}
\newcommand{\Ophtsqrtnbeta}{\Op^{\arght}(\tmpratetwo)}
\newcommand{\ophtsqrtnbeta}{\op^{\arght}(\tmpratetwo)}
\newcommand{\numun}{U_n}
%We now use the size of $\PkH$, so that this vanishes at the rate $\tmprate$.
Assumption \ifnum\jasa=1 5 \else \ref{assump:Cho1} \fi implies that $\inf_{\histlower \in \history} \PkH \ge \knot \tmprate$ for all $\histlower \in \history$ where $\knot > 0 $ is some constant. Note that $\tmprateinv \times n^{-1/2} = \tmpratetwo$. Multiplying $\tmprateinv$ to the numerator and denominator of the right-hand side of \cref{rcif_nj_0}, we obtain
\newcommand{\denomlong}{\tmprateinv \PkH + \Ophtsqrtnbeta}
\begin{align}
\frac{\PnkH\Nj\argt}{\PnkH} = &\frac{\tmprateinv \PkH\Nj\argt + \Ophtsqrtnbeta}{\denomlong}. \label{rcif_nj_1.1}
\end{align}
The denominator of the right-hand side of \cref{rcif_nj_1.1} is greater than or equal to $\knot + \opht$ because $\tmprateinv \PkH \ge \tmprateinv\knot\tmprate =\knot$ and the residual error is going to zero since by Assumption \ifnum\jasa=1 5\else \ref{assump:Cho1}\fi, $\beta>\frac{1}{2}$. Therefore, we have
\begin{equation}
    \begin{split}
\frac{\PnkH\Nj\argt}{\PnkH} = \frac{\tmprateinv\PkH\Nj\argt+\opht}{\tmprateinv\PkH+\opht}. \label{rcif_nj_1.2}
    \end{split}
\end{equation}
%Let $\numun = \tmprateinv \PkH\Nj\argt + \Ophtsqrtnbeta$.

We note that
\begin{equation}
\label{rcif_nj_2}
%\implies \frac{\numun}{\denomlong} &= \frac{\numun}{\tmprateinv \PkH} \bigg( 1 + \opht \bigg) \eqnspace
\frac{\tmprateinv \PkH + \opht}{\tmprateinv\PkH}
= 1+\frac{\opht}{\tmprateinv\PkH}
=1+\opht,
\end{equation}
where the last equality in \cref{rcif_nj_2} holds because $\tmprateinv\PkH \ge \knot > 0$.
Since $(1+\opht)^{-1}=1+\opht$, we can take the inverse of \cref{rcif_nj_2}.
\begin{align}
\implies \frac{\tmprateinv\PkH}{\tmprateinv \PkH + \opht} &= 1+\opht \notag \eqnspace
%divide each side by \tmprateinv\PkH
\implies \frac{1}{\tmprateinv \PkH + \opht} &= \frac{1 + \opht}{\tmprateinv \PkH}. \label{rcif_3}
\end{align}

To continue from \cref{rcif_nj_1.2}, multiply both sides of \cref{rcif_3} by $\tmprateinv \PkH\Nj\argt + \opht$.
\begin{align*}
%multiple each side by \tmprateinv \PkH\Nj\argt + \opht
&\implies \frac{\tmprateinv \PkH\Nj\argt + \opht}{\tmprateinv \PkH + \opht} = \frac{\tmprateinv \PkH\Nj\argt + \opht}{\tmprateinv \PkH} \bigg( 1 + \opht \bigg). %\notag \eqnspace
\end{align*}

We re-write the numerator as $\tmprateinv\PkH\Nj\argt=\tmpratetwo n^{\frac{1}{2}}\PkH\Nj\argt$. Since \\ $\sqrtn\PkH\Nj\argt = \Opht$ and $\tmpratetwo\to 0$, $\tmprateinv\PkH\Nj\argt$ is bounded. The denominator is bounded away from $0$ because $\tmprateinv\PkH \ge \knot>0$. This implies,
\begin{align*}
\frac{\PnkH\Nj\argt}{\PnkH} &=\frac{\tmprateinv \PkH\Nj\argt}{\tmprateinv \PkH} \bigg( 1 + \opht \bigg) + \opht %\label{residual}
\eqnspace
%can cancel out \tmprateinv now.
&=\frac{\PkH\Nj\argt}{\PkH} \bigg( 1 + \opht \bigg) + \opht \notag \eqnspace
%\frac{\PkH\Nj\argt}{\PkH} is bounded quantity, and little o_p times bounded quantity is still little o_p
&=\frac{\PkH\Nj\argt}{\PkH} + \opht \quad \text{(because $\frac{\PkH\Nj\argt}{\PkH}$ is bounded%, so $\frac{\PkH\Nj\argt}{\PkH} \times \opht = \opht$
).} 
%\label{lastline}
\end{align*}

Hence, $$\frac{\PnkH\Nj\argt}{\PnkH} = \frac{\PkH\Nj\argt}{\PkH} + \opht.$$ %and thus we have consistency of the kernel, or convergence for $\Nj\argt$.

Using similar arguments, we can show that the kernel version of $\atrisk$,  $\frac{1}{\truesurv\argtleft}\PnkH\atrisk\argt$ is consistent as well. Specifically, since $\truesurv\argt$ is bounded away from $0$, similar to \cref{rcif_nj_1.1}, we have
\begin{align*}
\frac{\frac{1}{\truesurv\argtleft}\PnkH\atrisk\argt}{\PnkH} = &\frac{\tmprateinv \frac{1}{\truesurv\argtleft}\PkH\atrisk\argt + \Ophtsqrtnbeta}{\denomlong}%\label{rcif_y_1.1} \eqnspace
&=\frac{\tmprateinv\frac{1}{\truesurv\argtleft}\PkH\atrisk\argt+\opht}{\tmprateinv\PkH+\opht}. %\label{rcif_y_1.2}
\end{align*}
Multiplying $\tmprateinv \frac{1}{\truesurv\argtleft}\PkH\atrisk\argt + \opht$ on the both sides of \cref{rcif_3} and simplifying the expressions, we have
\begin{align*}
%multiple each side by \tmprateinv \frac{1}{\truesurv\argtleft}\PkH\atrisk\argt + \opht
\implies \frac{\tmprateinv \frac{1}{\truesurv\argtleft}\PkH\atrisk\argt + \opht}{\tmprateinv \PkH + \opht} %&= \frac{\tmprateinv \frac{1}{\truesurv\argtleft}\PkH\atrisk\argt + \opht}{\tmprateinv \PkH} \bigg( 1 + \opht \bigg) \eqnspace
%&=\frac{\tmprateinv \frac{1}{\truesurv\argtleft}\PkH\atrisk\argt}{\tmprateinv \PkH} \bigg( 1 + \opht \bigg) + \opht \eqnspace
%%can cancel out \tmprateinv now.
%&=\frac{\frac{1}{\truesurv\argtleft}\PkH\atrisk\argt}{\PkH} \bigg( 1 + \opht \bigg) + \opht \eqnspace
%\frac{\frac{1}{\truesurv\argtleft}\PkH\atrisk\argt}{\PkH} is bounded quantity, and little o_p times bounded quantity is still little o_p
&=\frac{\frac{1}{\truesurv\argtleft}\PkH\atrisk\argt}{\PkH} + \opht.
\end{align*}
By Assumption \ifnum\jasa=1 5\else \ref{assump:CuiAssump1}\fi, $\PkH\atrisk\argt \ge \tauUB \PkH$. Then $\frac{\frac{1}{\truesurv\argtleft}\PkH\atrisk\argt}{\PkH} \ge \frac{\frac{1}{\truesurv\argtleft}\tauUB \PkH}{\PkH} = \frac{\tauUB}{\truesurv\argtleft}$. Since $\truesurv\argtleft$ is bounded away from $0$, we have $\frac{\frac{1}{\truesurv\argtleft}\PnkH\atrisk\argt}{\PnkH}$ is bounded away from $0$.
} % end of tmp space

\begin{comment}
% Bibliography
\bibliographystyle{apalike}
\bibliography{00_references/Overview, 00_references/00_references_manual}
%\end{document}
\end{comment}

\begin{comment}
%General Idea
Then, we can use Lemma 12.3 with continuous mapping theorem. Since $\Ah\argt$ is bounded below by $0$ and the map $\frac{1}{\Ah\argt}$ is smooth, we get $\frac{\truesurv\argtleft}{\atrisk
\argt}$. Because we have Hadamard differentiable functions, the map is continuous. We then want to use continuous mapping theorem, since each term is consistent so the continuous map is also consistent. With our assumptions, the continuous function of $\Bh\argt$ and $\Ah\argt$ gives the consistency of the random cumulative incidence forest.
\end{comment}

Let $\Ah\argu = \frac{1}{\truesurv\arguleft}\frac{\PnkH\atrisk\argu}{\PnkH}$ and $\Bh\argu = \frac{\PnkH\Nj\argu}{\PnkH}$. Let $\Bh_0\argt= \frac{\PkH \Nj\argt}{\PkH}$ and $\Ah_0\argt = \invtruesurv \frac{\PkH \atrisk \argt}{\PkH}$. Then our point of interest is $(B,A) = (\Bh_0\argt,\Ah_0\argt)$. The components of the second term of our estimator $\AJhat\argtgh$ in \cref{eq:rcif_est} can be mapped twice:
\begin{align*}
    (B,A) \mapsto (B,\frac{1}{A}) \mapsto \int_{[0,t]} \frac{1}{\Ah\argu}\,d\Bh \argu.
\end{align*}

Using Lemma 12.3 and Lemma 6.19 from \cite{kosorok_introduction_2008}, this composition map is Hadamard differentiable on a domain of type 
$\{ (B,A) :
\suph\int_{[0,\tau]} \lvert d\Bh\argt \rvert \leq M,\\
\inf_{\substack{t \in [0, \tau] \\ \histlower \in \history}} \lvert \Ah\argt \rvert \geq \epsilon \}$ 
for all $\histlower \in \history$ and for a given \( M < \infty \) and \( \epsilon > 0 \), at every point \((B,A)\) such that \( \frac{1}{\Ah\argt} \) has bounded variation. Thus, the map is continuous. Using Assumption \ifnum\jasa=1 9\else \ref{khbddtv}\fi, the total variation of the kernels is bounded. Since we have shown that $\Ah\argt$ and $\Bh\argt$ are consistent, and $\Ah\argt$ is bounded below by $0$ and the map $\frac{1}{\Ah\argt}$ is smooth, we apply the continuous mapping theorem so that the continuous map is also consistent.

Thus,
\begin{equation*}
\suptzerototauh \abs{\AJhat \argtgh - \cuminc_j \argtgh} \to 0
\end{equation*}
in probability as $n \to \infty$.

\begin{comment}
Then by Lemma 12.3, the derivative of the composition map is given by $$ (\alpha,\beta) \mapsto (\alpha, \frac{-\beta\truesurv^2}{\atrisk_0^2}) \mapsto \int_{[0,t]} \frac{\truesurv d\alpha}{\atrisk_0} - \int_{[0,t]} \frac{\beta \truesurv dN_0}{\atrisk_0^2}$$.

Then, from \cref{eq:rcif_distr} and using Theorem 2.8 in \cite{kosorok_introduction_2008}, we obtain \begin{align*}
    \sqrtn (\AJhat - \cuminc_j) \rightsquigarrow \int_{[0,(\cdot)]} \frac{\truesurv d\Gone}{\atrisk_0} - \int_{[0,(\cdot)]} \frac{\truesurv\Gtwo dN_0}{\atrisk_0^2}
\end{align*}
\end{comment}
\begin{comment}
Next, we obtain a bound for $\dLamnhHat$ vs. $\dLamnhTilde$ for $t \in (0,\tau]$. Using Riemann-Stieltjes integration by parts, we have
\begin{align*}
\intzerotot \truesurv\argsleft &\left[ \dLamnhHat\args - \dLamnhTilde\args \right] \eqnspace
&= \truesurv\argtleft \left[\LamnhHat\argt - \LamnhTilde\argt\right] - \truesurv(0) \left[ \LamnhHat(0) - \LamnhTilde(0)
\right] 
-\intzerotot \left[ \LamnhHat\args - \LamnhTilde\args \right] d\truesurv\args
\eqnspace
&= \truesurv\argtleft \left[\LamnhHat\argt - \LamnhTilde\argt\right] 
-\intzerotot \left[ \LamnhHat\args - \LamnhTilde\args \right] d\truesurv\args
\eqnspace
&\leq 2 \suptzerototauh \left| \LamnhHat\argt - \LamnhTilde\argt \right|
%this has the rate from the derivation from Lemma 12.3?
\end{align*}

For Lemma 12.3 and bounded kernels (also verifying that total variance is bounded), we have:
\begin{align*}
    \sqrt{n}
\begin{pmatrix}
\survnhhat\argt - \truehsurv\argt \\
\Pn \kernel\argH \Nj\argt - \PkH \Nj\argt\\
\Pn \kernel\argH \atrisk\argt - \PkH \atrisk \argt
\end{pmatrix}
= \sqrt{n}
\begin{pmatrix}
\survnhhat\argt - \truehsurv\argt \\
\PnP \knj\\
\PnP \ky
\end{pmatrix}
= \Opht %uniform over \histlower and t
\end{align*}

 Then,
\begin{align}
&\intzerotou \shat \dLamnHAT - \intzerotou \strue \dLamnTILDE = \intzerotou \shat\frac{\dLamnHATnum}{\dLamnHATdenom} - \intzerotou \strue \frac{\dLamnTILDEnum}{\dLamnTILDEdenom} \notag \eqnspace
&= \intzerotou \shat \dLamnHAT - \strue \dLamnTILDE + \left[ \strue \dLamnHAT - \strue \dLamnHAT \right] \notag \eqnspace
&= \intzerotou \left[ \shat - \strue \right] \dLamnHAT \\
&\quad \quad \quad + \intzerotou \tmpone - \intzerotou \strue 
\left\{\frac{\PnP \ky}{\dLamnHATdenom}\right\}\dLamnTILDE, \notag
\end{align}
because $\strue\dLamnHAT - \strue\dLamnTILDE$
\begin{align*}
%\strue\dLamnHAT &- \strue\dLamnTILDE \eqnspace
&= \strue\left[\frac{\dLamnHATnum-\dLamnTILDEnum+\dLamnTILDEnum}{\dLamnHATdenom} - \frac{\dLamnTILDEnum}{\dLamnTILDEdenom}\right] \eqnspace 
&= \strue\left[ 
\frac{\PnP \kdnj}{\Pn \ky} + \frac{\dLamnTILDEnum}{\dLamnHATdenom} - \frac{\dLamnTILDEnum}{\dLamnTILDEdenom}
\right]\eqnspace
&= \strue\left[ 
\frac{\PnP \kdnj}{\Pn \ky} - \frac{\dLamnTILDEnum}{\dLamnTILDEdenom} + \frac{\dLamnTILDEnum}{\dLamnHATdenom} 
\right]\eqnspace
&= \tmpone - \strue \left\{\frac{\dLamnTILDEnum}{\dLamnTILDEdenom} - \frac{\dLamnTILDEnum}{\dLamnHATdenom} 
\right\}\eqnspace
&= \tmpone - \strue 
\dLamnTILDEnum
\left\{\frac{\dLamnHATnum - \dLamnTILDEnum}{\dLamnTILDEdenom\dLamnHATdenom} 
\right\}\eqnspace
&= \tmpone - \strue 
\left\{\frac{\PnP \ky}{\dLamnHATdenom}\right\}\frac{\dLamnTILDEnum}{\dLamnTILDEdenom} \eqnspace
&= \tmpone - \strue 
\left\{\frac{\PnP \ky}{\dLamnHATdenom}\right\}\dLamnTILDE
\end{align*}

%next part
We know $\sqrtn \thetadiff = \widetilde{O}_p(1)$. Let $\fthat - \fttrue$ be a Hadamard differentiable function. We note that $\thetanHAT$ could be infinite dimensional.

Using Lemma 12.3 from Kosorok 2008 Springer, we can show that
\begin{align*}
\sqrtn \left(\fthat - \fttrue\right) &= \fdot \left( \sqrtn \thetadiff \right) + \op \eqnspace
&= \int^1_0 \fdotlong \left(\sqrtn \thetadiff\right)dt \eqnspace
&=\fdot \left( \sqrtn \thetadiff \right) + \int^1_0 \left[ \fdotlong - \fdot \right]dt \sqrtn \thetadiff
\end{align*}

In other words, $\fdotlong - \fdot \to 0$ within the bounded set $\thetadiff$ when $0<t<1$. We have convergence as $\sqrtn \thetadiff = \widetilde{O}_p(1) \implies \rn \norm{\thetanHAT-\truetheta} = \Opht$.
\end{comment}

\end{proof}
} %end command space

\begin{comment}
% Bibliography
\bibliographystyle{apalike}
\bibliography{00_references/Overview, 00_references/00_references_manual}
\end{document}
\end{comment}

} \fi
}


\newcommand{\PRELIMproofITR}{
\if1\prelim{
%The proof for Theorem \ref{thm:ITR_Consistency} is below.
%%%%%%%%%%%%%%%%%%%%%%%%%%%%%%%%%%%%%%%%%%%%%%%%%%%%%%%%%%%
%%%%%%%%%%%%%%%%%%%%%%%%%%%%%%%%%%%%%%%%%%%%%%%%%%%%%%%%%%%
%%%%%%%%%%%%%%%%%%%%%%%%%%%%%%%%%%%%%%%%%%%%%%%%%%%%%%%%%%%
% Consistency of ITR estimator
\newcommand{\ITRconsistency}{
\begin{theorem} (Consistency of the individualized treatment regime estimator) 
\label{thm:ITR_Consistency}
Assume Assumptions \ref{assump:CuiAssump1}-\ref{assump:CuiAssump2} hold. Let $\doptHAT$ denote the optimal ITR estimator that is built following Algorithm \ref{alg:Overview} and Assumptions \ref{assump:boundeddensity}-\ref{assump:boundedV0}. Assume Assumptions \ref{assump:Causal1}-\ref{assump:Causal3} hold and that the random forests are built for each treatment arm. Then, for %each $\phi = \phi^{\mu}, \phi^{\sigma,t_0}$ where 
$\phi^\mu = \begin{pmatrix}
    \phi^\mu_1 \\ \phi^\mu_2
\end{pmatrix},$ % and $\phi^{\sigma, t_0} = \begin{pmatrix} \phi^{\sigma, t_0}_1 \\ \phi^{\sigma, t_0}_2 \end{pmatrix},$
\[ \lim_{n \to \infty} \| \mathbb{V}(\doptTRUE(Z)) - \mathbb{V}(\doptHAT(Z)) \| \xrightarrow{p} 0\]
where $\|\cdot\|$ is defined as the $L_2$ norm.
\end{theorem}
}

\newcommand{\ITRconsistencyPROOF}{
\begin{proof}[Proof of Theorem \ifnum\jasa=1 1\else \ref{thm:ITR_Consistency}\fi]
%We know that, by properties of norm, $\left\| \begin{pmatrix} a_1 \\ a_2 \end{pmatrix} + \begin{pmatrix} b_1 \\ b_2 \end{pmatrix}\right\| \leq \left\| \begin{pmatrix} a_1 \\ a_2 \end{pmatrix} \right\| + \left\| \begin{pmatrix} b_1 \\ b_2 \end{pmatrix} \right\|$ and $\left\| c\begin{pmatrix} a_1 \\ a_2 \end{pmatrix} \right\| =  |c| \left\|\begin{pmatrix} a_1 \\ a_2 \end{pmatrix}\right\|$ where $c$ is a scalar. Also, by property of expectation, we know that $\|\mathbb{E}(\cdot)\| \leq \mathbb{E}(\|\cdot\|)$.

Note that
\begin{equation*}
\begin{split}
\left\| \mathbb{V}(\doptTRUE(Z)) - \mathbb{V}(\doptHAT(Z)) \right\| 
% line 1
&= \left\|\sum_{a \in \mathcal{A}} \mathbb{E}_Z\left[ \valueTRUE
\cdot \left( \mathbf{1}\{\doptTRUE(Z) = a\} - \mathbf{1}\{\doptHAT(Z) = a\} \right)
\right] \right\| \\
% line 2
&\leq \sum_{a \in \mathcal{A}} \left\|\mathbb{E}_Z\left[ \valueTRUE
\cdot \left( \mathbf{1}\{\doptTRUE(Z) = a\} - \mathbf{1}\{\doptHAT(Z) = a\} \right)
\right] \right\| \\
% line 3
&\leq \sum_{a \in \mathcal{A}} \mathbb{E}_Z \left\| \left[ \valueTRUE
\cdot \left( \mathbf{1}\{\doptTRUE(Z) = a\} - \mathbf{1}\{\doptHAT(Z) = a\} \right)
\right] \right\| \\
% line 4
&\leq \sum_{a \in \mathcal{A}} \mathbb{E}_Z \left[ \left\| \valueTRUE \right\|
\cdot \left| \mathbf{1}\{\doptTRUE(Z) = a\} - \mathbf{1}\{\doptHAT(Z) = a\} \right|
\right] \\
% line 5
&= \mathbb{E}_Z \left[ \sum_{a \in \mathcal{A}} \left( \mathbf{1} \{\doptTRUE(Z) \neq \doptHAT(Z)\} \left\|\valueTRUE\right\| \right) \right] \\
% line 6
&\leq \mathbb{E}_Z \left( \mathbf{1} \{\doptTRUE(Z) \neq \doptHAT(Z)\} \right) \sup_{a,Z} \left\| \valueTRUE \right\| \\
% line 7
&= \text{Pr}(\doptTRUE(Z) \neq \doptHAT(Z) | \doptHAT) \sup_{a,Z} \left\| \valueTRUE \right\|\\
% line 8
&\leq \sevenbound \text{Pr}(\doptTRUE(Z) \neq \doptHAT(Z) | \doptHAT).\\
\end{split} % split gives one number; align gives a number per line
\end{equation*} The third to last line holds because 1) only one term in the summation will have a non-zero contribution and 2) the supremum represents the ``worst case scenario" over all fixed combinations of $a$ and $z$ and thus does not depend on the random variables. The last line holds because we assumed $\sup_{a,Z} \left\| \valueTRUE \right\|$ is bounded by $\sevenbound$ via Assumption \ifnum\jasa=1 11\else \ref{assump:boundedV0} \fi. Thus, the norm of this distance, $\mathbb{V}(\doptTRUE) - \mathbb{V}(\doptHAT)$, goes to $0$ in probability if $\text{Pr}(\doptTRUE(Z) \neq \doptHAT(Z)|\hat{d}_\text{opt}) \xrightarrow{p} 0$.\\

We complete the proof by first showing
\begin{align}
    \label{eqn:step1}
        \sum_{a \in \mathcal{A}} \left| I\{ \valueoneTRUE - (1-\tol)\valueonemaxTRUE \ge 0 \} - I\{ \valueoneHAT - (1-\tol)\valueonemaxHAT \ge 0 \} \right| \rightarrow 0.
        \end{align}
    %(Note: we take the difference because both are random (one relies on Z and the other is an estimator)).
The proof of \cref{eqn:step1} will be presented shortly. Now, \cref{eqn:step1} implies 
\begin{align}
    \label{eqn:star}
        \argmax_{a \in \mathcal{A}: \valueoneHAT \ge (1-\tol)\valueonemaxHAT} \valuetwoHAT \xrightarrow{p} \argmax_{a \in \mathcal{A}: \valueoneTRUE \ge (1-\tol)\valueonemaxTRUE} \valuetwoHAT.
    \end{align} 
Theorem \ifnum \jasa=1 3 \else \ref{thm:RCIF_Consistency} \fi in our paper states that RCIFs are consistent. %replace <- with theorem when typed up.
Thus, by the Continuous Mapping Theorem, $\valuetwoHAT \xrightarrow{p} \valuetwoTRUE$. %$\valuetwoHAT$ is bounded.
Then,
\begin{align}
\label{eqn:step2}
        \sup_{a \in \mathcal{A}} \, \mathbb{E} \, \left| \valuetwoHAT - \valuetwoTRUE \right| \xrightarrow{p} 0,
    \end{align} 
which implies that the right hand side of \cref{eqn:star} converges to \\$\argmax_{a \in \mathcal{A}: \valueoneTRUE \ge (1-\tol)\valueonemaxTRUE} \valuetwoTRUE$. Thus, since $Z$ is random, $$\text{Pr}(\doptTRUE(Z) \neq \doptHAT(Z) | \doptHAT) \xrightarrow{p} 0.$$ This completes the proof of Theorem \ifnum\jasa=1 1\else \ref{thm:ITR_Consistency}\fi.
\end{proof}
}

\newcommand{\eqnsteponePROOF}{
%\textbf{Proof of \cref{eqn:step1}:}
\begin{proof}[Proof of \cref{eqn:step1}]
Let $A_0 = \valueoneTRUE - (1-\tol)\valueonemaxTRUE$
and $\widehat{A}_n = \valueoneHAT - (1-\tol)\valueonemaxHAT$. These are functions of random variable $Z$. Then, we can re-write \cref{eqn:step1} as 
$$\sum_{a \in \mathcal{A}} \left| I\{ A_0 \ge 0 \} - I \{ \widehat{A}_n \ge 0 \} \right| \rightarrow 0.$$
\begin{equation}
\label{eqn:terms12}
\begin{split}
    \left| I\{ \widehat{A}_n \ge 0 \} - I \{A_0 \ge 0 \} \right| 
    &= I\{ \widehat{A}_n < 0 \le A_0 \} + I \{A_0 < 0 \le \widehat{A}_n \} \\
    &= I\{ \widehat{A}_n - A_0 < -A_0 \leq 0 \} + I \{0 < -A_0 \le \widehat{A}_n - A_0 \} \\
    &= I\{ 0 \le A_0 < -(\widehat{A}_n - A_0) \} + I \{- (\widehat{A}_n - A_0) \le A_0 < 0 \}. \\
\end{split}
\end{equation}
In \cref{eqn:terms12}, the first equality holds because these are the only possible indicators such that their sum is equal to $1$. We get the 2nd equality by subtracting $A_0$ on both sides within the indicator functions. The third equality is due to multiplying by $-1$ and flipping the inequalities within the indicators.

Since we know $\valueoneHAT \xrightarrow{p} \valueoneTRUE$ and $\valueonemaxHAT \xrightarrow{p} \valueonemaxTRUE$ as shown in \cite{cho_multi-stage_2023}; and $\widehat{A}_n$ and $A_0$ are linear combinations of $\valueoneHAT$ and $\valueonemaxHAT$, and $\valueoneTRUE$ and $\valueonemaxTRUE$, respectively, then %Theorem \ref{thm:consistency_lincomb},
$\widehat{A}_n \xrightarrow{p} A_0$ because the linear combination of consistent estimators is consistent. This means
\begin{align}
\label{eqn:AnA0} 
\widehat{A}_n - A_0 \xrightarrow{p} 0.
\end{align}

By definition of convergence in probability, \cref{eqn:AnA0} implies $\exists\epsilon_n: Pr(|\widehat{A}_n - A_0|>\epsilon_n) \rightarrow 0 \text{ as } \epsilon_n \rightarrow 0$ for a fixed sequence $\epsilon_n$. Then, using Assumption \ifnum\jasa=1 10 \else \ref{assump:boundeddensity} \fi and \cref{eqn:AnA0}, we can bound \cref{eqn:terms12}. Starting with the first term of \cref{eqn:terms12},
\begin{equation*}
\label{eqn:term1}
\begin{split}
&I\{ 0 \le A_0 < -(\widehat{A}_n - A_0), |\widehat{A}_n-A_0|\le \epsilon_n \}  + I\{ 0 \le A_0 < -(\widehat{A}_n - A_0), |\widehat{A}_n-A_0| > \epsilon_n \} \le \\&I\{ 0 \le A_0 < \epsilon_n\} + I \{ |\widehat{A}_n - A_0| > \epsilon_n \}, \\    
\end{split}
\end{equation*}
 where $I\{ 0 \le A_0 < \epsilon_n\} \rightarrow 0$ by Assumption \ifnum\jasa=1 10 \else \ref{assump:boundeddensity}\fi (bounded density assumption) and
$I \{ |\widehat{A}_n - A_0| > \epsilon_n \} \rightarrow 0$ by \cref{eqn:AnA0} (and definition of consistency). Specifically, assuming Assumption \ifnum\jasa=1 10\else \ref{assump:boundeddensity}\fi,
$Pr(0\le A_0 \le \epsilon_n) = \int^{\epsilon_n}_{0} f_{A_0}(a_0)da_0 \approx f_{a_0}(0)\epsilon_n \le b\epsilon_n \rightarrow 0 $ as $\epsilon_n \rightarrow 0 $ for some bound $b$. The second term can be shown to be bounded similarly. 
\begin{equation*}
\begin{split}
&I \{- (\widehat{A}_n - A_0) \le A_0 < 0 \} = \\
&I \{- (\widehat{A}_n - A_0) \le A_0 < 0 , |\widehat{A}_n - A_0| \le \epsilon_n\} +
I \{- (\widehat{A}_n - A_0) \le A_0 < 0 , |\widehat{A}_n - A_0| > \epsilon_n\} \leq\\
&I(-\epsilon_n \leq A_0 < 0) + I(|\widehat{A}_n - A_0| > \epsilon_n).
\end{split}
\end{equation*} The first term for the inequality is because $-(\widehat{A}_n - A_0) < 0$, $\widehat{A}_n - A_0$ must be positive; thus, $-\epsilon_n \leq \widehat{A}_n - A_0 \leq \epsilon_n$.  The second term is because the indicator is a subset.

Because of Equation \cref{eqn:AnA0}, $Pr(|\widehat{A}_n - A_0| > \epsilon_n) \rightarrow 0$, and $Pr(-\epsilon_n \leq A_0 < 0) \rightarrow 0$ due to Assumption \ifnum\jasa=1 10\else \ref{assump:boundeddensity}\fi.
Specifically, $I(-\epsilon_n \leq A_0 < 0)  \implies Pr(-\epsilon_n \le A_0 < 0) \approx \epsilon_n f_{a_0}(0) < \epsilon_n b \rightarrow 0$ as $\epsilon_n \rightarrow 0$.
This completes the proof of \cref{eqn:step1}.
\end{proof}
}

\begin{comment}
MARCH 2024 UPDATE: WE CAN IGNORE BELOW I THINK.

We want to use the below equation. I think he said that, as long as the density of the quantity $\valueoneTRUE - (1-\tol)\valueonemaxTRUE$ is bounded then the whole thing converges to 0?
\begin{equation}
%\begin{align}
    \lim_{n \to \infty} E\left(\left|I( \valueoneHAT - (1-\tol)\valueonemaxHAT \ge 0) - I( \valueoneTRUE - (1-\tol)\valueonemaxTRUE \ge 0 )\right|\right) = 0 \\
%\end{align}
\end{equation}
We want to show that $\valueoneHAT - (1-\tol)\valueonemaxHAT -\xrightarrow{p} \valueoneTRUE - (1-\tol)\valueonemaxTRUE$
Let $X_n \xrightarrow{p} \mu$ and $Y_n \xrightarrow{p} \gamma $. Then, $Pr(|X_n - \mu| > \epsilon \text{ or } |Y_n - \gamma|>\epsilon) \leq Pr(|X_n - \mu|>\epsilon) + Pr(|Y_n - \gamma|>\epsilon)$. If $n \rightarrow \infty$, then both $Pr(|X_n - \mu|>\epsilon) \rightarrow 0$ and $Pr(|Y_n - \gamma|>\epsilon) \rightarrow 0$. We want to use this as the generic idea for $X_n$ being OS and $Y_n$ being CIF.
\end{comment}

\begin{comment}
MARCH 2024 UPDATE: DEFINITELY IGNORE BELOW. PROOF IS COMPLETE.

\subsubsection{OLD STUFF: Notes from before 2/20/24 meeting}

Sample size $n$ such that $\hat{d}_n$. Let 
\[
V(\hat{d}_n) = \begin{pmatrix}
V_1(\hat{d}_n) \\ V_2(\hat{d}_n)
\end{pmatrix}
\]

We want to show:
\[
\lim_{{n \to \infty}} \bigg( \frac{V_1(d_*) - V_1(\hat{d}_n)}{V_1(d_*)} - \tol \bigg)_+ \xrightarrow{p} 0
\]
\\

Our optimal $d_{**}$ satisfies the following two criteria:
\begin{enumerate}
    \item $E\left[ I(\frac{V_1(d_*|Z)-V_1(d_{**}|Z)}{V_1(d_*|Z)} \leq \tol ) \right] = 1$
    \item maximize $E\left[V_2(d_{**}|Z) \;\middle|\; I(\frac{V_1(d_*|Z)-V_1(d_{**}|Z)}{V_1(d_*|Z)} \leq \tol ) = 1\right]$
\end{enumerate}
\end{comment}
\ITRconsistencyPROOF
\eqnsteponePROOF
} \fi
}


\PRELIMproofRCIF
\fi
\fi
